\section{Proofs for Section~\ref{sec:constraints}}
\label{app:constraints}

This appendix proves the results of Section~\ref{sec:constraints}. We keep
the notation of Section~\ref{sec:constraints-setting}. Rational numbers,
vectors, and polynomials are encoded in binary, and $\bits(\cdot)$ denotes
encoding length. A \emph{rational circuit} is a pair of integer circuits
$(N,Q)$ over $\{0,1,+,-,\times\}$ with $Q>0$, denoting $N/Q$; sums,
differences, and products add a constant number of gates, division by a
nonzero $N_2/Q_2$ uses
\[
 \frac{N_1/Q_1}{N_2/Q_2}=\frac{N_1Q_2N_2}{Q_1N_2^2},
\]
and the sign of $N/Q$ is the sign of $N$ (Definition~\ref{def:models-circuits} and Lemma~\ref{lem:denominator-clearing}).
Comparing two rational circuits costs at most two $\PosSLP$ queries. For
their difference numerator $N$, equality holds exactly when both strict-sign
tests $N>0$ and $-N>0$ return false.

External results are used only through the following contracts. Each
contract is stated where it is used.
\begin{enumerate}[label=(E\arabic*)]
\item Convex value optimization (Lemma~\ref{lem:convex-value}): for an
explicit rational convex polynomial of fixed degree, a nonempty explicit
rational polyhedron contained in a supplied rational box, and a rational
$\varepsilon\in(0,1]$, an exactly feasible rational point with objective value
at most the minimum plus $\varepsilon$ is computed in ordinary time polynomial
in the input length and $\log(1/\varepsilon)$.
\item Separation (Lemma~\ref{lem:separation}) and the exact unconstrained
comparisons of Theorems~\ref{thm:exact-upper} and~\ref{thm:upper-monotone},
including the uniform Newton circuit of Lemma~\ref{lem:newton-circuit} and
the compilation Theorem~\ref{thm:posslp-closure}.
\item Linear programming with an arithmetic count independent of objective
and right-hand-side lengths
\citep[Theorem~6.6.3]{GroetschelLovaszSchrijver1988}, exact convex quadratic
programming \citep{KozlovTarasovKhachiyan1980}, exact separable quadratic
flows \citep[Theorem~20]{Vegh2016}, Clarkson-type sampling for violator spaces
\citep{GaertnerMatousekRuestSkovron2008}, integer feasibility with
integer-point separation oracles \citep{AriHildebrand2026,HildebrandGoess2024,Basu2023Integers},
and fixed-dimensional integer and mixed-integer linear feasibility
\citep{HildebrandKoeppe2013,DelPia2023}.
\end{enumerate}

\subsection{Geometry, radii, and supports}
\label{app:constraints-geometry}

Two elementary consequences of strong convexity are used repeatedly. For
completeness we derive them in the form needed; they are instances of
Lemma~\ref{lem:models-strong-convexity}.

\begin{lemma}[Radius and distance]
\label{lem:constraints-radius}
Let $\varphi$ be a twice differentiable function on $\R^N$ with
$\nabla^2\varphi\succeq\mu I$.
\begin{enumerate}[label=(\alph*)]
\item Let $x_0\in\R^N$, put
$A_0=\|\nabla\varphi(0)\|_1+|\varphi(x_0)-\varphi(0)|+1$ and
$R_0=2+2A_0/\mu$. Then $\varphi(x)>\varphi(x_0)$ whenever $\|x\|\ge R_0$. In
particular, every minimizer of $\varphi$ over a closed set containing $x_0$
has norm less than $R_0$.
\item If $\varphi$ is minimized over a closed convex set $K$ at $p$, then
$\varphi(x)-\varphi(p)\ge\frac\mu2\|x-p\|^2$ for every $x\in K$.
\end{enumerate}
If $\varphi$ is an explicit rational polynomial of fixed degree and $x_0$ is
rational, then $R_0$ is rational with bit length polynomial in
$\bits(\varphi)+\bits(\mu)+\bits(x_0)$.
\end{lemma}

\begin{proof}
(a) Write $\rho=\|x\|$ and $\Delta=\varphi(x_0)-\varphi(0)$. Strong convexity
and $|\nabla\varphi(0)^{\mathsf T}x|\le\|\nabla\varphi(0)\|_1\rho$ give
$\varphi(x)\ge\varphi(0)-\|\nabla\varphi(0)\|_1\rho+\frac\mu2\rho^2$. If
$\rho\ge R_0$, then $\frac\mu2\rho-\|\nabla\varphi(0)\|_1\ge A_0-\|\nabla\varphi(0)\|_1=|\Delta|+1$,
so $\varphi(x)-\varphi(x_0)\ge\rho(|\Delta|+1)-\Delta>0$ because $\rho\ge2$.
(b) $\varphi(x)\ge\varphi(p)+\nabla\varphi(p)^{\mathsf T}(x-p)+\frac\mu2\|x-p\|^2$,
and $\nabla\varphi(p)^{\mathsf T}(x-p)\ge0$ by optimality of $p$ on $K$.
The size statement holds because $\varphi(0)$, $\varphi(x_0)$, and
$\nabla\varphi(0)$ are fixed-degree evaluations.
\end{proof}

\begin{proof}[Proof of Lemma~\ref{lem:constraints-polyhedral}]
(a) Fix $x_0\in P$ and $R_1>\|T(x_0)\|/\mu$, and let
$K=P\cap\{x:\|x-x_0\|\le R_1\}$. The map $x\mapsto\Pi_K(x-T(x))$, where
$\Pi_K$ is Euclidean projection, is a continuous self-map of the nonempty
compact convex set $K$, so Brouwer's theorem gives a fixed point $\bar p$.
The projection inequality at $\bar p$ reads $T(\bar p)^{\mathsf T}(x-\bar p)\ge0$
for $x\in K$. If $\|\bar p-x_0\|=R_1$, then \eqref{eq:constraints-monotone}
gives
\[
 T(\bar p)^{\mathsf T}(\bar p-x_0)\ge T(x_0)^{\mathsf T}(\bar p-x_0)+\mu R_1^2
 \ge R_1(\mu R_1-\|T(x_0)\|)>0,
\]
contradicting the inequality at $x=x_0$. Hence $\|\bar p-x_0\|<R_1$. For any
$x\in P$ and small $\varepsilon>0$, the point $\bar p+\varepsilon(x-\bar p)$
lies in $K$, so $T(\bar p)^{\mathsf T}(x-\bar p)\ge0$; thus $\bar p$ solves
$\mathrm{VI}(T,P)$. If $p,p'$ are two solutions, adding
$T(p)^{\mathsf T}(p'-p)\ge0$ and $T(p')^{\mathsf T}(p-p')\ge0$ gives
$(T(p)-T(p'))^{\mathsf T}(p-p')\le0$, and \eqref{eq:constraints-monotone}
forces $p=p'$. For $T=\nabla f$, integrating
\eqref{eq:constraints-curvature} along segments gives
\eqref{eq:constraints-monotone}, and convexity shows that the solutions of
$\mathrm{VI}(\nabla f,P)$ are exactly the minimizers of $f$ on $P$.

(b) If $-T(x)=\sum_{i\in S(x)}\lambda_ia_i$ with $\lambda\ge0$, then for
$y\in P$ we get $T(x)^{\mathsf T}(y-x)=-\sum_i\lambda_i(a_i^{\mathsf T}y-b_i)\ge0$.
Conversely, let $x$ solve the variational inequality and let $d$ satisfy
$a_i^{\mathsf T}d\le0$ for $i\in S(x)$. The rows outside $S(x)$ have positive
slack at $x$, so $x+\varepsilon d\in P$ for small $\varepsilon>0$, and
$T(x)^{\mathsf T}d\ge0$. By Farkas' lemma, $-T(x)$ is a nonnegative
combination of $\{a_i:i\in S(x)\}$. Choose such a combination with the
smallest number of positive coefficients, supported on $J$. If the rows
$a_i$, $i\in J$, were dependent, a nonzero $\beta$ with
$\sum_{i\in J}\beta_ia_i=0$ could be chosen with some $\beta_i>0$; subtracting
$\vartheta\beta$ with $\vartheta=\min\{\lambda_i/\beta_i:\beta_i>0\}$ would
remove a coefficient while keeping the others nonnegative. Hence the rows in
$J$ are independent and $|J|\le\rank A$; if $T(x)=0$, then $J=\varnothing$.

(c) Let $|J|=s$. Gaussian elimination with a fixed pivoting rule selects $s$
columns $\mathcal B$ for which $A_{J,\mathcal B}$ is invertible; let
$\mathcal N$ be the other columns. Setting $x_{\mathcal N}=y$ and
$x_{\mathcal B}=A_{J,\mathcal B}^{-1}(b_J-A_{J,\mathcal N}y)$ parametrizes
$\{x:A_Jx=b_J\}$ as $\bar x+Zy$, with $Z_{\mathcal N}=I$ and
$Z_{\mathcal B}=-A_{J,\mathcal B}^{-1}A_{J,\mathcal N}$. Thus
$Z^{\mathsf T}Z=I+Z_{\mathcal B}^{\mathsf T}Z_{\mathcal B}\succeq I$. After the
rows of $A$ are scaled to integers, every entry of $\bar x$ and $Z$ is a ratio
of minors of the scaled matrix $[A\ b]$, so Hadamard's inequality bounds all
bit lengths by one polynomial in $L$, for every $J$. For $y,y'\in\R^{n-s}$,
\[
 (U_J(y)-U_J(y'))^{\mathsf T}(y-y')
 =\bigl(T(\bar x+Zy)-T(\bar x+Zy')\bigr)^{\mathsf T}Z(y-y')
 \ge\mu\|Z(y-y')\|^2\ge\mu\|y-y'\|^2 .
\]
Substituting an affine map into a polynomial of fixed degree creates
polynomially many monomials with polynomially bounded coefficients. For
$T=\nabla f$, the chain rule gives $\nabla g_J(y)=Z^{\mathsf T}\nabla f(\bar x+Zy)$
and $\nabla^2g_J=Z^{\mathsf T}\nabla^2fZ\succeq\mu Z^{\mathsf T}Z\succeq\mu I$.

(d) Since $A_Jp=b_J$ and $Z$ has full column rank, $p=\bar x+Zy_p$ for a
unique $y_p$. Because $A_J(\bar x+Zy)=b_J$ for all $y$, we have $A_JZ=0$, so
$Z^{\mathsf T}a_i=0$ for $i\in J$. As $T(p)$ lies in the span of these rows,
$U_J(y_p)=Z^{\mathsf T}T(p)=0$. By (a) applied to $U_J$ on all of
$\R^{n-s}$, the zero of $U_J$ is unique, so $y_p=y_J$. A support $J$ from (b)
satisfies $A_Jp=b_J$ and $T(p)=-A_J^{\mathsf T}\lambda$. A row basis $J$ of
$A_{S(p)}$ satisfies $A_Jp=b_J$, and
$T(p)\in-\operatorname{cone}\{a_i:i\in S(p)\}\subseteq\Span\{a_i:i\in J\}$.
\end{proof}

The next lemma is the deterministic test behind the support certificate of
Section~\ref{sec:constraints-general} and behind the violation primitive of
Theorem~\ref{thm:constraints-rank}. For a set $G$ of rows write
$P_G=\{x:a_i^{\mathsf T}x\le b_i\ (i\in G)\}$.

\begin{lemma}[Support test]
\label{lem:constraints-support-test}
Let $G$ be a set of rows with $P_G\ne\varnothing$, let $p_G$ be the solution of
$\mathrm{VI}(T,P_G)$, and let $J\subseteq G$ be a set of independent rows with
chart $(\bar x,Z)$, zero $y_J$ of $U_J$, and $p_J=\bar x+Zy_J$. With
$\lambda_J$ as in \eqref{eq:constraints-multipliers}, consider the tests
\begin{equation}\label{eq:constraints-support-test}
 b_i-a_i^{\mathsf T}p_J\ge0\ (i\in G),\qquad \lambda_J(y_J)\ge0,\qquad
 T(p_J)+A_J^{\mathsf T}\lambda_J(y_J)=0 .
\end{equation}
If they hold, then $p_J=p_G$. Conversely, every support $J\subseteq G$ of
$p_G$ given by Lemma~\ref{lem:constraints-polyhedral}(b) passes them. Each
test is a sign condition on an explicit polynomial of degree at most three in
$y$, with coefficient length polynomial in $L$, at the zero of $U_J$.
Hence the tests are decided in polynomial time with a $\PosSLP$ oracle, by
Theorem~\ref{thm:upper-monotone}, or by rational arithmetic if $|J|=n$.
\end{lemma}

\begin{proof}
Suppose \eqref{eq:constraints-support-test} holds. Then $p_J\in P_G$,
$a_i^{\mathsf T}p_J=b_i$ for $i\in J$, and $-T(p_J)=A_J^{\mathsf T}\lambda$ with
$\lambda=\lambda_J(y_J)\ge0$. For $x\in P_G$,
$T(p_J)^{\mathsf T}(x-p_J)=-\lambda^{\mathsf T}(A_Jx-b_J)\ge0$. Thus $p_J$
solves $\mathrm{VI}(T,P_G)$, and $p_J=p_G$ by uniqueness. Conversely, if $J$
is a support of $p_G$ with $T(p_G)=-A_J^{\mathsf T}\lambda$ and $\lambda\ge0$,
then $p_J=p_G$ by Lemma~\ref{lem:constraints-polyhedral}(d), and
$\lambda_J(y_J)=(A_JA_J^{\mathsf T})^{-1}A_JA_J^{\mathsf T}\lambda=\lambda$,
so all tests hold. The matrix $(A_JA_J^{\mathsf T})^{-1}$ is rational of
polynomial size because the rows of $J$ are independent; the tested
expressions are affine or cubic in $y$ after substitution.
\end{proof}

Guessing $J$ and running these tests with $G=\{1,\ldots,m\}$, followed by a
comparison of $h(p_J)$, proves the
$\NP^{\PosSLP}\cap\mathrm{coNP}^{\PosSLP}$ bound mentioned in
Remark~\ref{rem:constraints-unambiguous}.

\subsection{Unambiguous certificates}
\label{app:constraints-unambiguous}

\paragraph{A linear program with a circuit objective.}
We use the following linear-programming theorem, whose arithmetic count is
independent of the objective and right-hand-side lengths.

\begin{quote}
\emph{Contract (E3a)}
\citep[Theorem~6.6.3]{GroetschelLovaszSchrijver1988}. There is an algorithm
that, given rational $D_0$, $e_0$, and $c$ such that
$\{v:D_0v\le e_0\}$ is nonempty and bounded, returns a vertex minimizing
$c^{\mathsf T}v$ over this polytope. Its operations on numbers are additions,
subtractions, multiplications, divisions by nonzero numbers, and comparisons,
and their number is bounded by a polynomial in $\bits(D_0)$, independently of
$c$ and $e_0$. The integer roundings in its preprocessing of the objective are
carried out by comparisons of numbers whose range is bounded in advance.
\end{quote}

\begin{lemma}[Linear programming with a circuit objective]
\label{lem:constraints-circuit-lp}
Let $Q_0=\{v\in\R^n:D_0v\le e_0\}$ be a nonempty bounded polytope with
explicit rational data, and let $c$ be a vector of rational circuits. A vertex
of $Q_0$ minimizing $c^{\mathsf T}v$ can be computed, with its coordinates
printed as rationals in binary, in deterministic time polynomial in
$\bits(D_0,e_0)$ and the circuit size, with a $\PosSLP$ oracle.
\end{lemma}

\begin{proof}
Run the algorithm of contract (E3a) with every number held as a rational
circuit. Binary rational constants are built from $0$ and $1$ by doubling and
addition. Each arithmetic operation adds a constant number of shared gates,
and each comparison is answered exactly by $\PosSLP$. The simulation therefore
follows the execution of the algorithm on the rational input $c$, all
divisions are by nonzero numbers, and the output is an optimal vertex
$v_*$ represented by circuits. The number of gates and queries is polynomial.

To print $v_*$, decide for each row $i$ whether $(D_0)_iv_*=(e_0)_i$, by two
strict-sign queries. The rows that are tight at a vertex span $\R^n$:
otherwise a nonzero $d$ orthogonal to all of them gives
$v_*\pm\varepsilon d\in Q_0$ for small $\varepsilon>0$, because the other rows
have positive slack, and $v_*$ would not be extreme. Choose $n$ independent
tight rows by rational elimination on the explicit matrix $D_0$ and solve the
corresponding square system. Its unique solution is $v_*$, and Cramer's rule
bounds its bit length polynomially in $\bits(D_0,e_0)$. No circuit is
expanded.
\end{proof}

\paragraph{The normal-cone test.}
We use Lemma~\ref{lem:newton-circuit} in the following form. Let $U$ be an
explicit cubic map on $\R^d$, $d\ge1$, satisfying
\eqref{eq:constraints-monotone} with modulus $\mu$, with unique zero $y_*$,
and let $M$ be an integer, written in unary, that is at least the encoding
length of $(U,\mu)$. Part~(b) of that lemma, with $L'=M$, constructs in time
polynomial in $M$ and without oracle calls a shared rational circuit $\hat y$
such that every explicit polynomial $\eta$ of encoding length at most $M$
satisfies
\begin{equation}\label{eq:constraints-uniform-newton}
 \eta(y_*)=0\ \text{or}\ |\eta(y_*)|\ge g_M,
 \qquad\text{and}\qquad |\eta(\hat y)-\eta(y_*)|\le g_M/8,
\end{equation}
where $g_M=2^{-2^{e(M)}}$ has a circuit of $e(M)$ squarings of $1/2$. For
gradient maps, part~(a) of the lemma gives the same conclusion with a warm
start from convex value optimization.

\begin{lemma}[Normal-cone test]
\label{lem:constraints-cone-test}
Let $U$ and $y_*$ be as above, let $\mathsf C(y)\in\Q[y]^n$ be an explicit
vector of polynomials of degree at most three, and let
$Q_0=\{v:D_0v\le e_0\}$ be a nonempty bounded rational polytope. Whether
$\mathsf C(y_*)^{\mathsf T}v\ge0$ for all $v\in Q_0$ is decidable in
deterministic polynomial time with a $\PosSLP$ oracle.
\end{lemma}

\begin{proof}
Every vertex of $Q_0$ is the unique solution of $n$ independent rows of
$D_0v=e_0$. After scaling rows to integers, Cramer's rule and Hadamard's
inequality bound the coordinates of every vertex by a polynomial $\sigma$ in
the input length $\Lambda_0$ of $(U,\mu,\mathsf C,D_0,e_0)$. Hence there is a
polynomial $M=M(\Lambda_0)\ge\Lambda_0$, computed without knowing any vertex,
that bounds the encoding length of every vertex observable
$\eta_v(y)=v^{\mathsf T}\mathsf C(y)$. Construct $\hat y$ as in
\eqref{eq:constraints-uniform-newton}, form the rational circuit vector
$\hat c=\mathsf C(\hat y)$, and use Lemma~\ref{lem:constraints-circuit-lp} to
print a vertex $w$ minimizing $\hat c^{\mathsf T}v$ over $Q_0$. Finally decide
$\eta_w(y_*)\ge0$ exactly, by Theorem~\ref{thm:upper-monotone} or by testing
the sign of $\eta_w(\hat y)+g_M/2$; answer accordingly.

If $\mathsf C(y_*)^{\mathsf T}v\ge0$ on $Q_0$, then in particular
$\eta_w(y_*)\ge0$. Otherwise the linear function
$v\mapsto\mathsf C(y_*)^{\mathsf T}v$ is negative somewhere on the polytope,
hence at some vertex $v_1$. Then $\eta_{v_1}(y_*)\le-g_M$ and
$\hat c^{\mathsf T}v_1\le-7g_M/8$, so $\hat c^{\mathsf T}w\le-7g_M/8$. If
$\eta_w(y_*)\ge0$, then $\hat c^{\mathsf T}w\ge-g_M/8$, which is
impossible. Thus
$\eta_w(y_*)<0$ and the test answers correctly. Ties among vertex values cause
no difficulty.
\end{proof}

The linear program has a rational objective given by a circuit; it is never
run with the algebraic vector $\mathsf C(y_*)$. The only algebraic quantity
whose sign is decided is the explicit cubic $\eta_w$ at $y_*$, so the
observable theorems are applied only to explicit polynomials.

\begin{proof}[Proof of Theorem~\ref{thm:constraints-unambiguous}]
We describe two nondeterministic oracle machines, $\mathsf M$ for the
promise problem and $\overline{\mathsf M}$ for its complement. Both first
perform the following deterministic steps. If a certificate is supplied and is
invalid, $\mathsf M$ rejects and $\overline{\mathsf M}$ accepts. If
$P=\varnothing$, decided by linear programming, $\mathsf M$ rejects and
$\overline{\mathsf M}$ accepts. If $n=0$, everything is a rational
computation. Otherwise both machines guess $S\subseteq\{1,\ldots,m\}$, one bit
per row, and then deterministically:
\begin{enumerate}[label=(\arabic*)]
\item choose the row basis $J$ of $A_S$ obtained by scanning the rows of $S$
in increasing order, and its chart $(\bar x,Z)$; let $y_J$ be the zero of
$U_J$ and $p_S=\bar x+Zy_J$;
\item reject unless $b_i-a_i^{\mathsf T}p_S=0$ for all $i\in S$ and
$b_i-a_i^{\mathsf T}p_S>0$ for all $i\notin S$;
\item reject unless \eqref{eq:constraints-cone-test} holds, decided by
Lemma~\ref{lem:constraints-cone-test} with $\mathsf C(y)=T(\bar x+Zy)$ and
$Q_0=Q_S$ (if $|J|=n$, by ordinary linear programming with the rational
objective $T(\bar x)$);
\item $\mathsf M$ accepts if $h(p_S)\bowtie0$, and $\overline{\mathsf M}$
accepts otherwise.
\end{enumerate}
Each slack in step~(2) and the observable in step~(4) is an explicit
polynomial in $y$ at the zero of $U_J$, which satisfies
\eqref{eq:constraints-monotone} by Lemma~\ref{lem:constraints-polyhedral}(c),
so Theorem~\ref{thm:upper-monotone} decides it. All steps take polynomial time
with a $\PosSLP$ oracle.

\emph{Soundness.} Suppose $S$ passes steps (2) and (3). Then $p_S\in P$ and
$A_Sp_S=b_S$. For $x\in P$, the vector $x-p_S$ satisfies
$A_S(x-p_S)\le0$, and $v=(x-p_S)/\max\{1,\|x-p_S\|_\infty\}$ lies in $Q_S$;
hence $T(p_S)^{\mathsf T}(x-p_S)\ge0$. Thus $p_S$ solves $\mathrm{VI}(T,P)$
and equals $p$. Step~(2) now says $S=S(p)$.

\emph{Completeness.} For $S=S(p)$, Lemma~\ref{lem:constraints-polyhedral}(d)
gives $p_S=p$, so step~(2) holds by the definition of $S(p)$. By
Lemma~\ref{lem:constraints-polyhedral}(b), $-T(p)=\sum_{i\in S(p)}\lambda_ia_i$
with $\lambda\ge0$, so $T(p)^{\mathsf T}v=-\sum_i\lambda_ia_i^{\mathsf T}v\ge0$
for $v\in Q_S$, and step~(3) holds.

Consequently exactly one guess reaches step~(4), and it uses $p_S=p$. On a
yes-instance $\mathsf M$ has exactly one accepting computation and
$\overline{\mathsf M}$ has none; on a no-instance the roles are exchanged.
Inconsistent guessed equalities, redundant rows, zero rows, zero multipliers,
and lower-dimensional polyhedra are all handled by the exact mask in
step~(2).
\end{proof}


\subsection{Few nonlinear directions}
\label{app:constraints-nonlinear}

\begin{proof}[Proof of Lemma~\ref{lem:constraints-essential}]
(a) For $v\in\R^n$,
$D_v(f_3+f_4)=\sum_iv_i\,\partial_i(f_3+f_4)=\sum_\beta(Wv)_\beta x^\beta$.
Since $D_vf_3$ and $D_vf_4$ are homogeneous of degrees two and three, both
vanish identically exactly when $Wv=0$. Each monomial of $f_3+f_4$
contributes at most four monomials to the partial derivatives, so $W$ has at
most $4L$ nonzero rows, and its rank, a row basis $U$ of its row space, and
$V=U^{\mathsf T}(UU^{\mathsf T})^{-1}$ are computed by rational elimination
in polynomial time.

(b) For $x\in\R^n$, the vector $e=x-VUx$ lies in $\ker U=K_f$. The derivative
of $\vartheta\mapsto(f_3+f_4)(VUx+\vartheta e)$ is
$D_e(f_3+f_4)(VUx+\vartheta e)=0$, so $(f_3+f_4)(x)=(f_3+f_4)(VUx)=\theta(Ux)$.

(c) Let $x=Mx'+a$ with $M$ invertible and put $f'(x')=f(Mx'+a)$. Expanding
$f_4(y+a)=f_4(y)+D_af_4(y)+(\text{terms of degree}\le2)$ and
$f_3(y+a)=f_3(y)+(\text{terms of degree}\le2)$ shows that the homogeneous
parts of $f'$ of degrees four and three are $f_4(Mx')$ and
$f_3(Mx')+(D_af_4)(Mx')$. For $v'\in\R^n$ and $v=Mv'$, the derivatives in
direction $v'$ are $(D_vf_4)(Mx')$ and $(D_vf_3+D_vD_af_4)(Mx')$. If
$D_vf_4\equiv0$, then $D_vD_af_4=D_aD_vf_4\equiv0$. Hence $v'\in K_{f'}$ if and
only if $D_vf_4\equiv0$ and $D_vf_3\equiv0$, that is, $K_{f'}=M^{-1}K_f$.
\end{proof}

The next lemma is the elimination on an arbitrary affine face. Its point is
the uniform size bound: the same polynomial $T$ works for every set of rows,
so the unknown active face of the optimizer can be used in proofs without
being found.

\begin{lemma}[Elimination of the quadratic directions]
\label{lem:constraints-elimination}
Let $f$ satisfy \eqref{eq:constraints-curvature} with nonlinear dimension $k$
and $U,V,\theta$ as in Lemma~\ref{lem:constraints-essential}. Let $J$ be a set
of independent rows with chart $(\bar x,Z)$, $Z\in\Q^{n\times d}$, and let
$y_J$ be the unconstrained minimizer of $y\mapsto f(\bar x+Zy)$. Rational
linear algebra computes $s=\rank(UZ)\le k$, $\bar c\in\Q^n$, and
$D\in\Q^{n\times s}$ of full column rank such that, for
$\varphi(u)=f(\bar c+Du)$:
\begin{enumerate}[label=(\roman*)]
\item $\nabla^2\varphi(u)\succeq\mu D^{\mathsf T}D\succ0$ for all $u\in\R^s$;
for $s\ge1$, the rational number
\[
 \nu=\frac{\mu\det(D^{\mathsf T}D)}{\tr(D^{\mathsf T}D)^{s-1}}
\]
satisfies $\mu D^{\mathsf T}D\succeq\nu I$ (Lemma~\ref{lem:models-det-trace});
\item $\bar x+Zy_J=\bar c+Du_*$, where $u_*$ is the unique zero of
$\nabla\varphi$;
\item for every observable $h$, the coefficients of $\bar c$, $D$, $\varphi$,
and $\eta(u)=h(\bar c+Du)$, after clearing denominators, have bit length at most
$T(L)$, where $T$ is a polynomial with an absolute exponent that depends on
neither $J$ nor $k$.
\end{enumerate}
\end{lemma}

\begin{proof}
Let $Y=UZ\in\Q^{k\times d}$. From the reduced row echelon form of $Y$, let
$W_0\in\Q^{d\times(d-s)}$ be the standard basis of $\ker Y$ (one vector per
nonpivot column) and let $T_0\in\Q^{d\times s}$ consist of the unit vectors of
the pivot columns. In the coordinates ordered as pivot and nonpivot columns,
$[T_0\ W_0]$ is block upper triangular with identity diagonal blocks, hence
invertible, and $YT_0$ has rank $s$ because $Y[T_0\ W_0]=[YT_0\ 0]$.

Write $y=T_0u+W_0w$ and $F(u,w)=f(\bar x+ZT_0u+ZW_0w)$. Since $UZW_0=0$,
Lemma~\ref{lem:constraints-essential}(b) gives
\[
 F(u,w)=f_{\le2}(\bar x+ZT_0u+ZW_0w)+\theta(U\bar x+YT_0u).
\]
The second term does not depend on $w$, and the first has degree at most two.
With $Q_2=\nabla^2f_{\le2}$, a constant rational matrix,
\[
 F(u,w)=\tfrac12w^{\mathsf T}Hw+w^{\mathsf T}(Gu+a)+F(u,0),
\]
where
\[
 H=(ZW_0)^{\mathsf T}Q_2ZW_0,\qquad G=(ZW_0)^{\mathsf T}Q_2ZT_0,\qquad
 a=(ZW_0)^{\mathsf T}\nabla f_{\le2}(\bar x).
\]
The Hessian of $F$ in $w$ equals $(ZW_0)^{\mathsf T}\nabla^2f(x)ZW_0$, because
the term $\theta(\cdot)$ does not depend on $w$; by
\eqref{eq:constraints-curvature} it is at least
$\mu(ZW_0)^{\mathsf T}ZW_0\succ0$, since $ZW_0$ has full column rank. Thus
$H\succ0$, and for fixed $u$ the unique minimizer in $w$ is
$w(u)=-H^{-1}(Gu+a)$. Put $\bar c=\bar x-ZW_0H^{-1}a$ and
$D=ZT_0-ZW_0H^{-1}G$, so that $\bar x+ZT_0u+ZW_0w(u)=\bar c+Du$ and
$\varphi(u)=\min_wF(u,w)$.

Since $UZW_0=0$, we have $UD=YT_0$, which has rank $s$; so $D$ has full
column rank, and (i) follows from
$\nabla^2\varphi=D^{\mathsf T}\nabla^2f(\bar c+Du)D$. For (ii), write
$y_J=T_0u_J+W_0w_J$. Minimizing first over $w$ shows $w_J=w(u_J)$ and that
$u_J$ minimizes $\varphi$; by (i), $u_J$ is the unique zero $u_*$ of
$\nabla\varphi$. For (iii), the construction is a fixed sequence of rational
operations: the chart, the product $UZ$, an echelon form, the matrices $H$,
$G$, $a$ and $H^{-1}$, and the substitution of the affine map $\bar c+Du$ into
$f$ and $h$. Each step multiplies, inverts, or row-reduces matrices of
dimension at most $n\le L$ whose entries have polynomially bounded bit length;
by Cramer's rule and Hadamard's inequality each such step preserves a
polynomial bound, and the number of steps is fixed, so one polynomial bound
holds for every $J$. Substituting an affine map in $s\le n$ variables into a
monomial of degree at most four gives at most $(n+1)^4$ terms, with additive
growth of coefficient bit length. Clearing denominators by the product of all
denominators keeps the bit length polynomial. None of these bounds depends on
$k$ beyond $k\le n$.
\end{proof}

\begin{proof}[Proof of Theorem~\ref{thm:constraints-nonlinear}]
\emph{A uniform gap.} Assume $P\ne\varnothing$ and let $J$ be a row basis of
$A_{S(p)}$. By Lemma~\ref{lem:constraints-polyhedral}(d), $p=\bar x+Zy_J$
with $y_J$ the unconstrained minimizer of $f(\bar x+Zy)$, so
Lemma~\ref{lem:constraints-elimination} gives $p=\bar c+Du_*$ and
$h(p)=\eta(u_*)$. The formula
\[
 \exists u\in\R^s:\quad \partial_1\varphi(u)=0\wedge\cdots\wedge
 \partial_s\varphi(u)=0\wedge z-\eta(u)=0
\]
defines the singleton $\{h(p)\}$, because $u_*$ is the only real zero of
$\nabla\varphi$. After clearing denominators it has $s\le k$ quantified
variables and $s+1$ integer polynomials of degree at most four with
coefficients of bit size at most $T(L)$. Lemma~\ref{lem:separation}, with
$s$ quantified variables, degree $d=4$, and $\tau=T(L)$, gives
$|h(p)|\ge2^{-2T(L)4^{c_0s}}$ when $h(p)\ne0$. With $G(k)=2\cdot4^{c_0k}$
this yields
\begin{equation}\label{eq:constraints-nonlinear-gap}
 h(p)\ne0\quad\Longrightarrow\quad|h(p)|\ge\gamma:=2^{-\omega},\qquad
 \omega=G(k)\,T(L).
\end{equation}
The bound is linear in $T(L)$, with a factor that depends only on $k$. If
$s=0$, then $h(p)=\eta()$ is rational of bit length at most $T(L)$, and the
same bound holds. The bound \eqref{eq:constraints-nonlinear-gap} holds for every
observable of encoding length at most $L$, in particular for the slack
observables $b_i-a_i^{\mathsf T}x$. The set $S(p)$ is used only to prove it.

\emph{The algorithm.} Decide $P=\varnothing$ by linear programming; otherwise
obtain a rational $x_0\in P$ of polynomial bit length. Compute $k$, $U$, $V$
(Lemma~\ref{lem:constraints-essential}), the radius $R_0$ of
Lemma~\ref{lem:constraints-radius}(a) for $f$ and $x_0$, so $\|p\|<R_0$, and
$\omega=G(k)T(L)$. Let $C_h$ be the sum of the absolute values of the
coefficients of $h$, and put
\[
 B=1+\max\Bigl\{4nC_h(R_0+1)^3,\ \max_i\|a_i\|_1\Bigr\},\qquad
 \delta=\min\Bigl\{1,\frac{\gamma}{8B}\Bigr\},\qquad
 \varepsilon=\min\Bigl\{1,\frac{\mu\delta^2}2\Bigr\}.
\]
On the ball $\|x\|\le R_0+1$, each partial derivative of a monomial
$cx^\alpha$ of degree at most four is bounded by $4|c|(R_0+1)^3$, so $B$
bounds $\|\nabla h\|$ and the gradients of all slack observables there.
Contract (E1), applied to $f$ on $P\cap\{x:\|x\|_\infty\le R_0\}$, returns a
rational $\hat x\in P$ with $f(\hat x)\le f(p)+\varepsilon$. Its running time
is polynomial in $L$ and $\log(1/\varepsilon)=O(\omega+\bits(B)+\bits(\mu))$,
hence at most $F_1(k)L^{C_1}$. By Lemma~\ref{lem:constraints-radius}(b),
$\|\hat x-p\|\le\delta\le1$, so the segment from $p$ to $\hat x$ lies in the
ball and $|h(\hat x)-h(p)|\le B\delta\le\gamma/8$. Return $0$ if
$|h(\hat x)|<\gamma/2$ and the sign of $h(\hat x)$ otherwise. If $h(p)=0$, then
$|h(\hat x)|\le\gamma/8$; if $h(p)\ne0$, then $|h(\hat x)|\ge7\gamma/8$ and the
signs agree. The rational evaluations take time polynomial in $L$ and $\omega$.

The same point $\hat x$ decides every slack, because $B$ also bounds their
gradients; this yields $S(p)$. Finally compute a row basis $J$ of $A_{S(p)}$
and apply Lemma~\ref{lem:constraints-elimination} to output $\bar c$, $D$, and
$\varphi$. All steps take time $F(k)L^{C}$ with an absolute $C$: the only
superpolynomial quantity is the number $\omega$ of accuracy bits, which enters
polynomially.
\end{proof}

The algorithm prints $\omega=2\cdot4^{c_0k}T(L)$ bits. Since
$\log(1/\varepsilon)=O(\omega+\poly(L))$ and every step is polynomial in $L$
and $\omega$, one can take $F(k)=2^{O(k)}$.

\subsection{Exact constrained Newton steps}
\label{app:constraints-newton}

\begin{proof}[Proof of Theorem~\ref{thm:constraints-transfer}]
Let $(f,P)\in\mathcal C$, let $h$ be an observable of degree at most $D$, and
let $L$ be the total input length. Decide $P\ne\varnothing$ by linear
programming and assume $n\ge1$.

\emph{Constants.} By linear programming compute a rational point
$x_{00}\in P$, and let $R_0$ be the radius of
Lemma~\ref{lem:constraints-radius}(a) for $f$ and $x_{00}$, so that
$\|p\|<R_0$. Put $\varrho=R_0+1$,
$\Omega=\{x:\|x\|_\infty\le\varrho\}$, and let $C_f$ and $C_h$ be the sums of
the absolute values of the coefficients of $f$ and $h$, increased to at least
one. Define
\[
 \Lambda=n^2D^3C_f\varrho^D,\qquad \kappa=\max\{1,\Lambda/(2\mu)\},\qquad
 B_h=1+nDC_h\varrho^D .
\]
On $\Omega$, every third partial derivative of $f$ is bounded by
$D^3C_f\varrho^D$. Hence each Hessian entry is
$\sqrt n\,D^3C_f\varrho^D$-Lipschitz there, and the Frobenius norm shows that
$\Lambda$ is a Lipschitz constant of $\nabla^2f$ on $\Omega$ in operator norm.
Similarly $\|\nabla h\|\le B_h$ on $\Omega$. All these numbers are rational
of polynomial bit length.

\emph{Warm start and convergence.} Contract (E1), applied to $f$ on
$P\cap\{x:\|x\|_\infty\le R_0\}$ with accuracy
$\varepsilon_0=\min\{1,\mu/(8\kappa^2)\}$, returns a rational $x_0\in P$ with
$f(x_0)\le f(p)+\varepsilon_0$, so $\|x_0-p\|\le1/(2\kappa)\le1/2$ by
Lemma~\ref{lem:constraints-radius}(b). Let $x_{j+1}=N(x_j)$. Every iterate
lies in $P$. If $\|x_j-p\|\le1/2$, then $x_j$ and $p$ lie in the convex set
$\Omega$, and Lemma~\ref{lem:constraints-newton-step} gives
$\|x_{j+1}-p\|\le\kappa\|x_j-p\|^2$. With $e_j=\kappa\|x_j-p\|$ we get
$e_0\le1/2$ and $e_{j+1}\le e_j^2$, so by induction
\begin{equation}\label{eq:constraints-newton-rate}
 \|x_j-p\|\le\kappa^{-1}2^{-2^j}\le\tfrac12\qquad(j\ge0).
\end{equation}

\emph{Separation.} The formula in the free variable $z$
\[
 \exists x\in\R^n\ \exists\lambda\in\R^m:\
 Ax\le b,\ \lambda\ge0,\ \nabla f(x)+A^{\mathsf T}\lambda=0,\
 \lambda_i(b_i-a_i^{\mathsf T}x)=0\ (i\le m),\ z=h(x)
\]
defines $\{h(p)\}$. Indeed, Lemma~\ref{lem:constraints-polyhedral}(b) at $p$
gives a multiplier supported on $S(p)$. Conversely, a solution $(x,\lambda)$
has $x\in P$ and $-\nabla f(x)\in\operatorname{cone}\{a_i:i\in S(x)\}$ by
complementarity, so $x$ solves $\mathrm{VI}(\nabla f,P)$ and equals $p$. The
formula has $n+m$ quantified variables and polynomially many integer
polynomials of degree at most $D$ with coefficients of polynomial bit length.
Lemma~\ref{lem:separation}, with $s=n+m$ quantified variables, degree at
most $\max\{D,2\}$, and coefficient bit size $\tau$ polynomial in $L$, gives
$h(p)=0$ or $|h(p)|\ge2^{-2\tau\max\{D,2\}^{c_0(n+m)}}$. Hence there is a
polynomial $e_D$, depending only on $D$, with $h(p)=0$ or
$|h(p)|\ge\gamma:=2^{-2^{e_D(L)}}$; enlarge $e_D$ so that also
$2^{e_D(L)}\ge\log_2B_h+3$. Then \eqref{eq:constraints-newton-rate} with
$j_0=e_D(L)+1$ gives
\[
 |h(x_{j_0})-h(p)|\le B_h\,2^{-2\cdot2^{e_D(L)}}\le2^{-2^{e_D(L)}-3}=\gamma/8 .
\]

\emph{Circuit implementation.} The point $x_0$ is printed in binary and turned
into rational circuits. For $j<j_0$, the entries of $H_j=\nabla^2f(x_j)$ and
$c_j=\nabla f(x_j)-H_jx_j$ are computed from the circuits of $x_j$ with
polynomially many gates, since the degree is fixed. The solver is then run on
$(H_j,c_j)$ with every number held as a rational circuit and every comparison
answered exactly by $\PosSLP$. Because the answers are exact, the simulated
execution is the execution of the solver on the rational data $(H_j,c_j)$:
every division is by a nonzero number and the output is $x_{j+1}=N(x_j)$. The
number of operations per step is polynomial in $L$ and independent of the bit
lengths of $(H_j,c_j)$, so after $j_0$ steps all circuits have polynomial size.
Finally build $\gamma$ by $e_D(L)$ squarings of $1/2$, let $\hat\alpha=h(x_{j_0})$,
and ask for the sign of the numerator of
\[
 \hat\alpha-\gamma/2,\qquad \hat\alpha+\gamma/2,\qquad\text{or}\qquad
 \gamma^2/4-\hat\alpha^2,
\]
which is positive exactly when $h(p)>0$, $h(p)\ge0$, or $h(p)=0$,
respectively; negations give the other relations. A zero value gives
$|\hat\alpha|\le\gamma/8$, and a nonzero value gives
$|\hat\alpha|\ge7\gamma/8$ with the same sign.

All work is polynomial in $L$, so the predicate is in $\mathrm P^{\PosSLP}$.
Applying the procedure to the slacks $b_i-a_i^{\mathsf T}x$ gives $S(p)$. If
$f(p)<\tau$, then for $h=f-\tau$ we have $h(p)\le-\gamma$, hence
$f(x_{j_0})-\tau\le-7\gamma/8<0$, and $x_{j_0}\in P$ is the claimed circuit
point. The comparisons inside the solver depend on earlier answers, so the
procedure is a Turing reduction.
\end{proof}

\subsubsection{Boxes}
\label{app:constraints-boxes}

\begin{lemma}[Inactive truncation]
\label{lem:constraints-truncation}
Let $f$ satisfy \eqref{eq:constraints-curvature}, let $P\ne\varnothing$ have
minimizer $p$, and let $q\in P$ be rational. Put
$R_1=1+2\|\nabla f(q)\|_1/\mu$ and
$P'=P\cap\{x:|x_i-q_i|\le R_1\ (i\le n)\}$. Then $|p_i-q_i|<R_1$ for all
$i$, so $p$ is the minimizer of $f$ on $P'$ and none of the added bounds is
active at $p$. The bit length of $R_1$ is polynomial in $L+\bits(q)$.
\end{lemma}

\begin{proof}
Since $f(p)\le f(q)$ and
$f(p)\ge f(q)-\|\nabla f(q)\|\,\|p-q\|+\frac\mu2\|p-q\|^2$, we get
$\|p-q\|\le2\|\nabla f(q)\|/\mu\le2\|\nabla f(q)\|_1/\mu<R_1$.
\end{proof}

For a box, $q$ is obtained by clipping $0$ to the finite bounds; for a flow
polyhedron, by linear programming. In both cases $P'$ is again a box or a flow
polyhedron, now with finite bounds, and the active bounds of $P'$ at $p$ are
the active original bounds.

\begin{lemma}[Admissible vector from the comparison matrix]
\label{lem:constraints-admissible}
Let $H$ be symmetric positive definite with $M(H)\succ0$. Put
$d=M(H)^{-1}\mathbf 1$ and $v=\frac12(H+M(H))d$. Then $d>0$, $v\ge\mathbf 1$,
and $H_{SS}^{-1}v_S>0$ entrywise for every nonempty $S\subseteq\{1,\ldots,n\}$.
If all off-diagonal entries of $H$ are nonpositive, then $v=\mathbf 1$.
\end{lemma}

\begin{proof}
\emph{Nonnegative inverse.} Let $B_0$ be positive definite with nonpositive
off-diagonal entries, and let $B_0u=w\ge0$. Write $u=u_+-u_-$ with
$u_\pm\ge0$ of disjoint supports. Then
$0\le w^{\mathsf T}u_-=u_-^{\mathsf T}B_0u_+-u_-^{\mathsf T}B_0u_-$. The first
term involves only off-diagonal entries, so it is nonpositive, and the second is
positive unless $u_-=0$. Hence $u\ge0$. Applied to $B_0=M(H)$ and $w=\mathbf 1$,
this gives $d\ge0$; if $d_i=0$, then $(M(H)d)_i=\sum_{j\ne i}M(H)_{ij}d_j\le0$,
contradicting $(M(H)d)_i=1$. So $d>0$.

\emph{Lower bound on $v$.} The identity $M(H)d=\mathbf 1$ reads
$H_{ii}d_i=1+\sum_{j\ne i}|H_{ij}|d_j$. Therefore
\[
 v_i=H_{ii}d_i+\sum_{j\ne i,\,H_{ij}<0}H_{ij}d_j
    =1+\sum_{j\ne i,\,H_{ij}>0}H_{ij}d_j\ge1 .
\]
If $H$ has nonpositive off-diagonal entries, the last sum is empty and
$v=\mathbf 1$.

\emph{Principal systems.} Fix a nonempty $S$ and consider the Jacobi map
$\mathcal J$ on $\R^S$ given by
$(\mathcal Jx)_i=\bigl(v_i-\sum_{j\in S\setminus\{i\}}H_{ij}x_j\bigr)/H_{ii}$.
Its fixed points are the solutions of $H_{SS}x=v_S$. For $x\in[0,d_S]$, the
numerator is smallest when $x_j=d_j$ for $H_{ij}>0$ and $x_j=0$ otherwise,
and then equals $1+\sum_{j\notin S,\,H_{ij}>0}H_{ij}d_j\ge1$. It is largest
when $x_j=d_j$ for $H_{ij}<0$ and $x_j=0$ otherwise, and then equals
$H_{ii}d_i+\sum_{j\notin S,\,H_{ij}<0}H_{ij}d_j\le H_{ii}d_i$. Hence
$\mathcal J$ maps $[0,d_S]$ into $[0,d_S]$, with every coordinate at least
$1/H_{ii}>0$. For $x,x'\in\R^S$, with $\|x\|_d=\max_i|x_i|/d_i$,
\[
 \frac{|(\mathcal Jx-\mathcal Jx')_i|}{d_i}
 \le\frac{\sum_{j\ne i}|H_{ij}|d_j}{H_{ii}d_i}\,\|x-x'\|_d
 =\Bigl(1-\frac1{H_{ii}d_i}\Bigr)\|x-x'\|_d ,
\]
so $\mathcal J$ is a contraction. Its unique fixed point lies in $[0,d_S]$ and
has positive coordinates; since $H_{SS}$ is invertible, the fixed point is
$H_{SS}^{-1}v_S$.
\end{proof}

The iteration $\mathcal J$ appears only in the proof. An algorithm computes
$M(H)$ by sign tests, $d$ by a positive definite solve, and $v$ by matrix
arithmetic.

\begin{lemma}[Parametric box quadratic programming]
\label{lem:constraints-box-qp}
Let $H\in\Q^{n\times n}$ be symmetric positive definite, $u\in\Q^n$ with
$u>0$, $c\in\Q^n$, and let $v\in\Q^n$, $v>0$, satisfy $H_{SS}^{-1}v_S\ge0$
for every nonempty $S$. The procedure below returns the minimizer of
$\frac12x^{\mathsf T}Hx+c^{\mathsf T}x$ over $[0,u]$ after at most $2n$
pivots. It uses $O(n^4)$ arithmetic operations and $O(n^2)$ comparisons, and
it divides only by pivots of positive definite matrices and by numbers it has
tested to be positive.
\end{lemma}

\emph{Procedure.} Maintain a partition $\{1,\ldots,n\}=\mathcal L\cup\mathcal
F\cup\mathcal U$ and a parameter $\theta\ge0$; initially $\mathcal L$ is
everything and $\theta=\max\{0,\max_i(-c_i/v_i)\}$. In each round compute,
using $LDL^{\mathsf T}$ factorization of $H_{\mathcal F\mathcal F}$,
\[
 \alpha_{\mathcal F}=H_{\mathcal F\mathcal F}^{-1}(c_{\mathcal F}+H_{\mathcal F\mathcal U}u_{\mathcal U}),\qquad
 \beta_{\mathcal F}=H_{\mathcal F\mathcal F}^{-1}v_{\mathcal F},
\]
and for $i\notin\mathcal F$,
$\alpha_i=c_i+H_{i\mathcal U}u_{\mathcal U}-H_{i\mathcal F}\alpha_{\mathcal F}$
and $\beta_i=v_i-H_{i\mathcal F}\beta_{\mathcal F}$. The \emph{event values}
are $-\alpha_i/\beta_i$ for $i\in\mathcal L$ with $\beta_i>0$ and
$-(u_i+\alpha_i)/\beta_i$ for $i\in\mathcal F$ with $\beta_i>0$. Let
$\theta'$ be the maximum of $0$ and all event values. If $\theta'=0$, return
the point $x$ with $x_{\mathcal L}=0$, $x_{\mathcal U}=u_{\mathcal U}$, and
$x_{\mathcal F}=-\alpha_{\mathcal F}$. Otherwise move the smallest index whose
event value equals $\theta'$ from $\mathcal L$ to $\mathcal F$ or from
$\mathcal F$ to $\mathcal U$, set $\theta=\theta'$, and start a new round.

\begin{proof}
For a partition and a parameter $\theta$, let $x(\theta)$ be the point with
$x_{\mathcal L}=0$, $x_{\mathcal U}=u_{\mathcal U}$, and
$x_{\mathcal F}=-\alpha_{\mathcal F}-\theta\beta_{\mathcal F}$. It solves the
free equations $(Hx+c+\theta v)_{\mathcal F}=0$, and for $i\notin\mathcal F$
the gradient of $q_\theta(x)=\frac12x^{\mathsf T}Hx+(c+\theta v)^{\mathsf T}x$
is $\partial_i(\theta)=\alpha_i+\theta\beta_i$. We prove that each round
starts with the following invariant: $0\le x(\theta)\le u$,
$\partial_i(\theta)\ge0$ for $i\in\mathcal L$, and $\partial_i(\theta)\le0$ for
$i\in\mathcal U$. These are the Karush--Kuhn--Tucker conditions for
minimizing $q_\theta$ over $[0,u]$. Initially $x=0$ and
$\partial_i=c_i+\theta v_i\ge0$ by the choice of $\theta$.

\emph{Signs of the slopes.} Admissibility gives $\beta_{\mathcal F}\ge0$. For
$i\notin\mathcal F$ let $S=\mathcal F\cup\{i\}$ and $w=H_{SS}^{-1}v_S\ge0$.
Eliminating $w_{\mathcal F}$ from $H_{SS}w=v_S$ gives
$(H_{ii}-H_{i\mathcal F}H_{\mathcal F\mathcal F}^{-1}H_{\mathcal Fi})\,w_i=\beta_i$.
The Schur complement on the left is positive, so $\beta_i\ge0$; for
$\mathcal F=\varnothing$, $\beta_i=v_i>0$.

\emph{One round.} As $\theta$ decreases with the partition fixed, the free
coordinates $-\alpha_i-\theta\beta_i$ do not decrease, and the gradients
$\alpha_i+\theta\beta_i$ outside $\mathcal F$ do not increase. Hence the
invariant persists until a gradient in $\mathcal L$ reaches zero or a free
coordinate reaches its upper bound. These are exactly the event values, and
each of them is at most the current $\theta$, because the invariant holds at
the current $\theta$; coordinates with zero slope never produce events. So the
invariant holds on $[\theta',\theta]$. If $\theta'=0$, it holds at $\theta=0$,
where $q_0$ is the original objective; the conditions are sufficient for
optimality because the problem is convex, and the minimizer is unique because
$H\succ0$.

\emph{A pivot.} If $\theta'>0$ and the selected index $i$ is in $\mathcal L$,
then $\partial_i(\theta')=0$ and $x_i=0$, so the current point satisfies the
free equations for $\mathcal F\cup\{i\}$. These equations have a unique
solution because $H_{\mathcal F\cup\{i\},\mathcal F\cup\{i\}}\succ0$, so the
new partition represents the same point at $\theta'$, and the invariant holds.
If $i\in\mathcal F$ with $x_i(\theta')=u_i$, the current point satisfies the
free equations for $\mathcal F\setminus\{i\}$ with $x_i=u_i$, and
$\partial_i(\theta')=0\le0$; again the invariant holds. Ties are processed one
index at a time.

An index moves only from $\mathcal L$ to $\mathcal F$ and from $\mathcal F$ to
$\mathcal U$, so there are at most $2n$ pivots and at most $2n+1$ rounds. Each
round uses $O(n^3)$ arithmetic operations; the $LDL^{\mathsf T}$ factorization
of a positive definite matrix needs no comparisons and has positive pivots.
Each round needs $O(n)$ comparisons to test $\beta_i>0$ and to find the
maximum and the smallest maximizing index.
\end{proof}

\begin{proof}[Proof of Corollary~\ref{cor:constraints-boxes}]
If some $\ell_i>u_i$, the box is empty. Substitute every fixed coordinate
$\ell_i=u_i$. If every coordinate is fixed, then $P=\{\ell\}$ and $p=\ell$ is
rational: evaluate $h(p)$ exactly, report every bound as active, and stop;
neither the transfer theorem nor the procedure of
Lemma~\ref{lem:constraints-box-qp} is called on an empty set of variables.
Otherwise the Hessian of the reduced objective is a principal submatrix of
$\nabla^2f$, so each of (i)--(iii) is inherited: principal submatrices of a
comparison matrix are comparison matrices, and forests and sign patterns
restrict. Apply Lemma~\ref{lem:constraints-truncation} to obtain a box with
finite bounds $\ell<u$ and the same minimizer; its added bounds are inactive.
Observables are substituted in the same way, and the active original bounds
are decided by the slack observables.

Conditions (i) and (ii) imply (iii). Under (ii), $M(H)=H$. Under (i), root each
tree of the forest, set $\sigma_{\mathrm{root}}=1$, and for an edge from a
parent $i$ to a child $j$ set $\sigma_j=-\operatorname{sgn}(H_{ij})\sigma_i$
if $H_{ij}\ne0$ and $\sigma_j=\sigma_i$ otherwise. With
$\Sigma=\diag(\sigma)$, every off-diagonal entry of $\Sigma H\Sigma$ is
$-|H_{ij}|$, so $M(H)=\Sigma H\Sigma\succ0$.

The Taylor-QP solver for the box translates $x=\ell+x'$, which replaces $c$ by
$c+H\ell$ and the box by $[0,u-\ell]$, computes $M(H)$ with $n^2$ sign tests,
$d$ and $v$ as in Lemma~\ref{lem:constraints-admissible} with $O(n^3)$
operations, and runs the procedure of Lemma~\ref{lem:constraints-box-qp}.
Under (iii) at the current point, Lemma~\ref{lem:constraints-admissible}
makes $v$ admissible, so the solver is exact. It uses $O(n^4)$ operations and
comparisons, independently of the bit lengths of $H$ and $c$, with a control
computable in polynomial time. It is therefore an arithmetic Taylor-QP solver,
and Theorem~\ref{thm:constraints-transfer} applies. Condition~(i) can be
checked from the coefficients after the substitution: the reduced box is
full-dimensional, and a polynomial that vanishes on a full-dimensional box
vanishes identically.
\end{proof}

\subsubsection{Separable flows}
\label{app:constraints-flows}

We use the quadratic case of V\'egh's strongly polynomial algorithm through
the following contract.

\begin{quote}
\emph{Contract (E3c)} \citep[Theorem~20 and Section~6.1]{Vegh2016}. There is
an algorithm that, given a directed graph with $m'\ge1$ arcs and no isolated
nodes, rational node demands, finite rational lower and upper arc capacities,
and arc costs $a_es^2+c_es$ with rational $a_e\ge0$ and $c_e$, decides whether
a feasible flow exists and, if so, returns an optimal flow. Its operations on
numbers are additions, subtractions, multiplications, divisions, and
comparisons; their number is $O(m'^4\log(m'+2))$ and does not depend on the
numerical data. On rational data the returned flow is the exact optimum.
\end{quote}

An isolated node carries no arc, so its demand must be zero in every feasible
flow. Isolated nodes are therefore deleted beforehand, with $O(\abs V)$
comparisons, and a nonzero demand at one of them shows that no feasible flow
exists.

Implementations of minimum-cost flow algorithms often make the network
strongly connected by artificial arcs of large linear cost. In the arithmetic
model, any such cost is computed from the data by the algorithm's own
operations. The following lemma gives, in addition, an explicit valid cost
that is a rational function of the data, so that this step never inspects the
bit length of a circuit.

\begin{lemma}[Explicit artificial costs]
\label{lem:constraints-artificial}
Consider a network with $N_0$ nodes, demands summing to zero, arcs $e$ with
capacity intervals $[0,U_e]$ (possibly $U_e=\infty$), and convex differentiable
costs $C_e$, and suppose that it has an optimal flow $s^*$. Let $C_0\ge1$ satisfy
$|C_e'(s^*_e)|\le C_0$ for all arcs. Add a node $\omega$ with zero demand and,
for each node $v$, uncapacitated arcs $(v,\omega)$ and $(\omega,v)$ with
linear cost $M_0s$, where $M_0>(N_0-1)C_0$. Then the optimal flows of the
augmented network are exactly the optimal flows of the original network,
extended by zero on the new arcs.
\end{lemma}

\begin{proof}
With node potentials $\pi$, the reduced cost of an arc $e=(i,j)$ at a flow
$s$ is $r_e=C_e'(s_e)-\pi_i+\pi_j$. A feasible flow is optimal if and only if
some potential satisfies, for every arc with $U_e>0$, $r_e\ge0$ when $s_e=0$,
$r_e=0$ when $0<s_e<U_e$, and $r_e\le0$ when $s_e=U_e$; an arc with $U_e=0$
imposes no condition, because the multipliers of its two bounds absorb any
reduced cost. This is the Karush--Kuhn--Tucker system of a convex problem
with linear constraints. Such potentials exist at $s^*$. In the
residual graph of $s^*$, which has an arc $i\to j$ of length $C_e'(s^*_e)$ when
$s^*_e<U_e$ and an arc $j\to i$ of length $-C_e'(s^*_e)$ when $s^*_e>0$, these
conditions say that $-\pi$ satisfies the triangle inequality, so there is no
negative cycle. Hence shortest-path distances $\delta(v)$ from an extra source
joined to every node by an arc of length zero are well defined, simple
shortest paths have at most $N_0-1$ residual arcs, and
$-(N_0-1)C_0\le\delta(v)\le0$. The potential $\pi=-\delta$ satisfies the
optimality conditions and takes values in $[0,(N_0-1)C_0]$. Set
$\pi_\omega=0$. The new arcs then have reduced costs $M_0-\pi_v>0$ and
$M_0+\pi_v>0$ at zero flow, so $s^*$ extended by zero is optimal in the
augmented network.

Let $s'$ be any feasible augmented flow and let $r$ be the reduced costs at
$s^*$ extended by zero. Convexity gives
$\mathrm{cost}(s')\ge\mathrm{cost}(s^*)+\sum_a c'_a(s'_a-s^*_a)$, where
$c'_a$ is $C_a'(s^*_a)$ or $M_0$. Since $s'$ and $s^*$ have the same demands,
$\sum_a(\pi_i-\pi_j)(s'_a-s^*_a)=0$, so the sum equals
$\sum_ar_a(s'_a-s^*_a)$. Each term is nonnegative: on an arc with $U_a>0$ by
the sign conditions, and on an arc with $U_a=0$ because
$s'_a=s^*_a=0$. Each new arc contributes $r_as'_a$ with $r_a>0$. Hence an
augmented optimum carries no flow on the new arcs, and its restriction is
optimal for the original network.
\end{proof}

For quadratic costs on finite capacity intervals, a valid $C_0$ is
$1+\sum_e\bigl(|C_e'(\text{lower bound})|+|C_e'(\text{upper bound})|\bigr)$,
because $C_e'$ is monotone. It is computed from circuit data with two sign
tests per arc.

\begin{proof}[Proof of Corollary~\ref{cor:constraints-flows}]
Decide feasibility by linear programming and obtain a rational feasible flow
$q$. Substitute every fixed arc ($\ell_e=u_e$) and adjust the demands. If no
arc remains, the feasible flow is the rational point $p$ given by the fixed
values: evaluate $h(p)$ exactly, report both bounds of every arc as active,
and stop. Otherwise delete the isolated nodes; their adjusted demands are
zero because $P\ne\varnothing$. Apply Lemma~\ref{lem:constraints-truncation}
with $q$ to obtain finite capacities with the same minimizer and inactive
added bounds; below, $\ell$ and $u$ denote these finite bounds. At a feasible
$x$, the Taylor model is, up to a constant, the separable quadratic flow
problem in the flow variable $y$ with arc costs $a_ey_e^2+c_ey_e$, where
$a_e=\frac12f_e''(x_e)>0$ and $c_e=f_e'(x_e)-f_e''(x_e)x_e$ are rational
circuits computed from those of $x$ with a constant number of gates per arc.
The shift $s=y-\ell$ turns the capacity intervals into $[0,u_e-\ell_e]$,
replaces the demands $b$ by $b-A_G\ell$, computed once by rational
arithmetic, and replaces $c_e$ by $c_e+2a_e\ell_e$, with a constant number of
gates per arc; it puts every subproblem in the form of
Lemma~\ref{lem:constraints-artificial}.
The algorithm of contract (E3c), run on the reduced graph, which has
$m'\ge1$ arcs and no isolated nodes, simulated on rational circuits with
$\PosSLP$ comparisons and, if needed, with the artificial cost of
Lemma~\ref{lem:constraints-artificial}, is an arithmetic Taylor-QP solver, and
Theorem~\ref{thm:constraints-transfer} applies. For $D\le4$, the promise
$f_e''\ge\mu$ is checked as follows: a quadratic $\alpha s^2+\beta s+\gamma_0$
is nonnegative on $\R$ exactly when $\alpha>0$ and $4\alpha\gamma_0\ge\beta^2$,
or $\alpha=\beta=0\le\gamma_0$.
\end{proof}

\subsubsection{One sign instance}
\label{app:constraints-single-sign}

We use Theorem~\ref{thm:posslp-closure}(b): if a deterministic oracle Turing
machine $\mathsf M$, whose oracle queries are encodings of integer circuits
over $\{0,1,+,-,\times\}$, halts within $\pi_0(L)$ steps on every input of
length $L$ and for all oracle answers, for a polynomial $\pi_0$, then there is
a polynomial-time map $x\mapsto\Phi_x$ to integer circuits with $\Phi_x>0$ if
and only if $\mathsf M$ with a $\PosSLP$ oracle accepts $x$.

\begin{proof}[Proof of Corollary~\ref{cor:structured-single-sign}]
Let $\mathsf M_0$ be the decision procedure of the proof of
Theorem~\ref{thm:constraints-transfer} for $(h,\bowtie)$, with the
preprocessing of Corollary~\ref{cor:constraints-boxes} or
Corollary~\ref{cor:constraints-flows} when it applies. Its queries are
numerators of rational circuits that it has written, hence encodings of
integer circuits. On every promised input of length $L$ its running time is at
most $\pi_1(L)$ for a fixed polynomial $\pi_1$: the ordinary steps (linear
programming, the warm start, the constants) are polynomial in bit time, and
the number of Newton steps, of solver operations per step, and of gates per
operation are bounded by polynomials in $L$ that do not depend on the expanded
values of any circuit. Each query has length at most the running time.

Let $\mathsf M$ check that its input is a well-formed instance, run
$\mathsf M_0$ for at most $\pi_1(L)$ steps, and reject if $\mathsf M_0$ has not
halted. Then $\mathsf M$ halts within a polynomial number of steps on every
input and agrees with $\mathsf M_0$ on promised inputs. On inputs that violate
a promise, a division by a number that happens to be zero only produces a
circuit with a zero denominator, which is still a syntactically valid query; the
clock bounds the computation. Invalid supplied certificates and empty feasible
sets are decided by ordinary steps before any query. The adaptive form of
Theorem~\ref{thm:posslp-closure}(b) now gives the required circuit $\Phi$. For a
polynomial-size Boolean combination of predicates, let one machine run the
corresponding procedures one after the other and evaluate the combination; it
satisfies the same hypotheses.
\end{proof}

The proof uses only that the decision procedure is deterministic and has a
worst-case polynomial bound on time and query length. A procedure with these
properties and a time bound $F(\kappa)L^C$ compiles into one circuit in time
polynomial in $F(\kappa)L^C$, which is an FPT many-one reduction. The Las Vegas
algorithm of Theorem~\ref{thm:constraints-rank} has only an expected time
bound, and fixing its random bits does not yield a correct deterministic
procedure; the guessed certificate of Theorem~\ref{thm:constraints-unambiguous}
is likewise not removed by compilation.

\subsection{Constraint rank}
\label{app:constraints-rank}

A \emph{violator space} \citep{GaertnerMatousekRuestSkovron2008} is a pair
$(H,V)$ of a finite set $H$ and a map $V:2^H\to2^H$ such that
$G\cap V(G)=\varnothing$ for all $G\subseteq H$ (consistency) and
$V(G)=V(F)$ whenever $F\subseteq G\subseteq H$ and $G\cap V(F)=\varnothing$
(locality). A \emph{basis} of $G$ is an inclusion-minimal $B\subseteq G$ with
$V(B)=V(G)$, and the \emph{combinatorial dimension} is the largest size of a
basis of any $G\subseteq H$.

\begin{lemma}[A violator space of rank $r$]
\label{lem:constraints-violator}
Let $P\ne\varnothing$, let $H=\{1,\ldots,m\}$ index the rows of $A$, and for
$G\subseteq H$ let $p_G$ be the solution of $\mathrm{VI}(T,P_G)$. Then
$V(G)=\{i\in H:a_i^{\mathsf T}p_G>b_i\}$ defines a violator space. Every basis
consists of linearly independent rows, so the combinatorial dimension is at
most $r=\rank A$. If $B$ is a basis of $H$, then $p_B=p$.
\end{lemma}

\begin{proof}
Every $P_G$ contains $P$, so $p_G$ exists and is unique by
Lemma~\ref{lem:constraints-polyhedral}(a). Consistency holds because
$p_G\in P_G$. For locality, $G\cap V(F)=\varnothing$ says that $p_F\in P_G$.
Since $P_G\subseteq P_F$, the inequality $T(p_F)^{\mathsf T}(x-p_F)\ge0$ holds
for all $x\in P_G$, so $p_F=p_G$ and $V(F)=V(G)$. No potential function or
ordering of values is used.

Let $B$ be a basis of some $G$. By Lemma~\ref{lem:constraints-polyhedral}(b),
applied to $\mathrm{VI}(T,P_B)$, there is a set $J\subseteq B$ of independent
rows with $a_i^{\mathsf T}p_B=b_i$ for $i\in J$ and
$-T(p_B)=A_J^{\mathsf T}\lambda$, $\lambda\ge0$. For $x\in P_J$,
$T(p_B)^{\mathsf T}(x-p_B)=-\lambda^{\mathsf T}(A_Jx-b_J)\ge0$, and
$p_B\in P_B\subseteq P_J$. Hence $p_J=p_B$, so $V(J)=V(B)=V(G)$, and
minimality of $B$ forces $J=B$. Finally, $V(H)=\varnothing$ because $p_H=p\in P$;
so a basis $B$ of $H$ has $V(B)=\varnothing$, that is, $p_B\in P$. Its
variational inequality holds on $P_B\supseteq P$, so $p_B=p$.
\end{proof}

The bound concerns every basis of every subset. A small active support at the
final solution alone would not bound the combinatorial dimension.

\begin{quote}
\emph{Contract (E3d)} \citep[Theorem~27 and Sections~4.2--4.5]{GaertnerMatousekRuestSkovron2008}.
Let $(H,V)$ be a violator space with $|H|=m$ and combinatorial dimension at
most a given $\delta\ge1$. Clarkson's two randomized algorithms, as formulated
for violator spaces, return a basis of $H$. They access $V$ only through the
\emph{violation test}, which decides $i\in V(G)$ for given $G\subseteq H$ and
$i\in H\setminus G$, and they perform an expected $O(\delta m+\delta^{O(\delta)})$
violation tests. In every test $|G|\le\delta$: the tests ask whether an element
violates a basis returned by a recursive call, or, inside the routine that
finds a basis of a sample of $O(\delta^2)$ elements by examining its subsets
$B$ with $|B|\le\delta$, whether an element violates $B$ or whether $b\in B$
violates $B\setminus\{b\}$. The remaining operations are random sampling of a
prescribed number of elements from a set or from a multiset given by integer
multiplicities, doubling of multiplicities, and set bookkeeping. In each call
of the reweighting algorithm the multiplicities are doubled $O(\delta\log m)$
times.
\end{quote}

\begin{lemma}[Violation primitive]
\label{lem:constraints-rank-primitive}
In the setting of Lemma~\ref{lem:constraints-violator}, given $G$ with
$|G|\le r$ and $i\notin G$, one decides $i\in V(G)$ in deterministic time
$2^rL^{C_0}$ with a $\PosSLP$ oracle, where $C_0$ is an absolute constant.
\end{lemma}

\begin{proof}
Enumerate the subsets $J\subseteq G$ with independent rows in a fixed order;
there are at most $2^r$. For each, compute the chart and $U_J$ and run the
tests \eqref{eq:constraints-support-test} with this $G$. By
Lemma~\ref{lem:constraints-support-test}, a passing $J$ gives $p_J=p_G$, and
some $J$ passes, namely a support of $p_G$, which lies in $G$. At the first
passing $J$, decide $a_i^{\mathsf T}p_J-b_i>0$, an affine observable at the
zero of $U_J$. Each $J$ costs polynomially many tests, each a $\PosSLP$
instance of polynomial length (Theorem~\ref{thm:upper-monotone}).
\end{proof}

\begin{proof}[Proof of Theorem~\ref{thm:constraints-rank}]
Decide $P=\varnothing$ by linear programming and compute $r$. If $r=0$, then
all rows are zero and feasibility gives $P=\R^n$; the solution is the zero of
$T$, and Theorem~\ref{thm:upper-monotone} decides the predicate. Otherwise run
Clarkson's algorithms of contract (E3d) with $\delta=r$ on the violator space
of Lemma~\ref{lem:constraints-violator}, implementing every violation test by
Lemma~\ref{lem:constraints-rank-primitive}. By contract (E3d), every test has a
first argument of size at most $r$, so the primitive applies. The returned
basis $B$ satisfies $p_B=p$ by Lemma~\ref{lem:constraints-violator}.
Enumerating the independent subsets of $B$ and running the support test with
$G=B$ yields $J$ and $(\bar x,Z,U_J)$ with $p=\bar x+Zy_J$. The predicate
$h(p)\bowtie0$ is then decided by Theorem~\ref{thm:upper-monotone} applied to
$U_J$ and $h(\bar x+Zy)$, or by rational arithmetic if $|J|=n$.

\emph{Bit model of the sampling.} Each element of $H$ is a row index, so
repeated inequalities are distinct elements. Multiplicities are stored as
binary integers. Since each call of the reweighting algorithm doubles
multiplicities $O(r\log m)$ times, every multiplicity, and their sum, has
$O(r\log m)$ bits. To draw one labeled copy from a multiset with total weight
$W$, draw a uniform integer in $\{1,\ldots,W\}$ by rejection from
$\lceil\log_2W\rceil$ random bits, which needs an expected constant number of
trials, and locate it by prefix sums; for sampling without replacement,
decrement the copy count temporarily. Before the basis routine, duplicate
copies are identified with their row index, since multiplicity affects
sampling probabilities but not the polyhedron $P_G$. Between consecutive
violation tests the algorithms perform a constant number of sampling rounds and
set operations on at most $m$ elements. Hence the expected bit time outside
the tests is polynomial in $L$ times one plus the expected number of tests.

\emph{Expected time and correctness.} The expected number of tests is
$O(rm+r^{O(r)})$ and each costs $2^rL^{C_0}$, with $m\le L$. The total expected
time is therefore $(r+1)^{O(r)}L^{C}$ with an absolute $C$. Every query has
length polynomial in $L$, independently of the random choices. The algorithms
stop only with a basis of $H$, and every test is answered exactly, so every
returned answer is correct; randomness affects only the running time.
\end{proof}

\subsection{Candidate lists with unconstrained fibers}
\label{app:constraints-lists}

We use the following integer feasibility theorem, which queries its oracle
only at integer points.

\begin{theorem}[Integer feasibility with integer-point separation;
\citealp{AriHildebrand2026,HildebrandGoess2024,Basu2023Integers}]
\label{thm:constraints-integer-query}
There is a deterministic algorithm with the following property. Its input is
an integer $t\ge1$ and an integer $R\ge2$, and it has access to an oracle for
a closed convex set $S\subseteq[-R,R]^t$ that is queried only at integer
points. On a query $z\in\Z^t$ the oracle either asserts $z\in S$, or returns
$(\pi,\beta)\in\Q^t\times\Q$ with $\pi^{\mathsf T}x\le\beta$ for all $x\in S$
and $\pi^{\mathsf T}z>\beta$. The algorithm either returns a point of
$S\cap\Z^t$ or correctly reports $S\cap\Z^t=\varnothing$. If the oracle's
answers on queries of bit length $\ell$ have length and computation time at
most $\Phi(\ell)$, its running time is
\[
 2^{O(t\log(t+1))}\poly\bigl(t,\ 1+\bits(R),\ \Phi(c_0t(1+\bits(R)))\bigr)
\]
for an absolute constant $c_0$ and a polynomial of absolute degree. Queries
outside $[-R,R]^t$ may be answered by a violated box inequality.
\end{theorem}

This is \citet[Definition~3.4 and Theorem~3.5]{AriHildebrand2026}, who
attribute it to \citet[Theorem~9 and Appendix~B]{HildebrandGoess2024} and
\citet[Theorem~5.7 and Remark~5.11]{Basu2023Integers}. No inradius or
full-dimensionality assumption on $S$ is required, empty and one-point sets
are allowed, and the queries are made in the original integer coordinates.
The algorithm is deterministic. Basu's remarks also contain the idea of
optimizing without bisection on the objective by keeping visited candidates,
which our lists use in a form that never needs objective values.

For the initial integer point we use deterministic integer linear feasibility
in fixed dimension: for an explicit rational polyhedron $Q\subseteq\R^t$, one
can decide $Q\cap\Z^t=\varnothing$ or return a point of $Q\cap\Z^t$ in time
$2^{O(t\log(t+1))}L^{O(1)}$, with output of bit length $2^{O(t)}L^{O(1)}$ and
absolute exponents \citep[Theorem~1.1, specialized to linear constraints and a
constant objective]{HildebrandKoeppe2013}.

\begin{lemma}[Gradient cuts for unconstrained fibers]
\label{lem:constraints-gradient-cut}
Let $P=Q\times\R^n$ and $g(z)=\min_yf(z,y)$ for $z\in\R^t$. Then $g$ is
differentiable and $\mu$-strongly convex on $\R^t$, with
$\nabla g(z)=\nabla_zf(z,y(z))$, where $y(z)$ is the unique fiber minimizer.
For an integer $R\ge1$ and $z\in\Z^t\cap[-R,R]^t$, a rational $c_z$ with
$\|c_z-\nabla g(z)\|\le\mu/4$, and hence satisfying
\eqref{eq:constraints-integer-cut}, is computed in ordinary time
$(L+\bits(R)+\bits(z))^{C_1}$ for an absolute constant $C_1$.
\end{lemma}

\begin{proof}
Each fiber $y\mapsto f(z,y)$ has Hessian at least $\mu I$, so $y(z)$ is
unique, and the implicit function theorem applied to $\nabla_yf(z,y)=0$ makes
$y(\cdot)$ smooth. Then $g(z)=f(z,y(z))$ and
$\nabla g(z)=\nabla_zf(z,y(z))+Dy(z)^{\mathsf T}\nabla_yf(z,y(z))=\nabla_zf(z,y(z))$.
For $z_1,z_2$ with fiber minimizers $y_1,y_2$ and $\vartheta\in[0,1]$,
strong convexity of $f$ gives
$g(\vartheta z_1+(1-\vartheta)z_2)\le\vartheta g(z_1)+(1-\vartheta)g(z_2)
-\frac\mu2\vartheta(1-\vartheta)\|z_1-z_2\|^2$, after dropping the
$y$-part of the squared distance. If $c_z$ satisfies the stated accuracy and
$w\ne z$ is an integer point, then
\[
 g(w)-g(z)\ge c_z^{\mathsf T}(w-z)-\tfrac\mu4\|w-z\|+\tfrac\mu2\|w-z\|^2
 \ge c_z^{\mathsf T}(w-z)+\tfrac\mu4\|w-z\|^2,
\]
because $\|w-z\|\ge1$.

To compute $c_z$, let $m_0=t+n$, $C_f=\max\{1,\sum_\alpha|f_\alpha|\}$, and
\[
 D_1=4m_0C_f(1+R)^3,\quad S_1=1+R+D_1/\mu,\quad K_1=1+12m_0C_fS_1^2,\quad
 \eta=\min\Bigl\{1,\frac{\mu}{4tK_1}\Bigr\}.
\]
Each partial derivative of $f$ at $(z,0)$ is at most $4C_f(1+R)^3$ in absolute
value, so $\|\nabla_yf(z,0)\|\le D_1$, and strong monotonicity of
$\nabla_yf(z,\cdot)$ gives $\|y(z)\|\le D_1/\mu$. On the set of points
$(z,y)$ with $\|y-y(z)\|\le1$, every coordinate is at most $S_1$ in absolute
value, every second partial derivative of $f$ is at most $12C_fS_1^2$, and the
block $\nabla^2_{zy}f$ has operator norm less than $K_1$. If $n=0$, put
$c_z=\nabla f(z)$. Otherwise contract (E1), applied to $y\mapsto f(z,y)$ on the
box $\|y\|_\infty\le D_1/\mu+1$, returns a rational $\hat y$ with
$f(z,\hat y)\le f(z,y(z))+\min\{1,\mu\eta^2/2\}$, hence $\|\hat y-y(z)\|\le\eta$.
Put $c_z=\nabla_zf(z,\hat y)$. For each $i\le t$, the mean value theorem along
the segment from $y(z)$ to $\hat y$ gives
$|(c_z)_i-\partial_ig(z)|\le K_1\eta\le\mu/(4t)$, so
$\|c_z-\nabla g(z)\|\le\sqrt t\,\mu/(4t)\le\mu/4$. All numbers involved have
bit length polynomial in $L+\bits(R)+\bits(z)$, and the accuracy requested in
(E1) has polynomially many bits.
\end{proof}

\begin{proof}[Proof of Theorem~\ref{thm:constraints-candidates}]
If $t=0$, the only possible block is the empty vector, which is feasible
exactly when $Q=\R^0$ satisfies its constraints. Let $t\ge1$. Use integer
linear feasibility to certify $Q\cap\Z^t=\varnothing$ or to find
$z_0\in Q\cap\Z^t$. Apply Lemma~\ref{lem:constraints-radius}(a) to $f$ on
$\R^{t+n}$ and the point $(z_0,0)$; every optimal solution has
$f\le f(z_0,0)$, so its norm is less than $R_0$. Let $R=\lceil R_0\rceil+1$
and $P_0=Q\cap[-R,R]^t$, which contains $z_0$ and every optimal integer block.
Since $\bits(z_0)\le2^{O(t)}L^{O(1)}$, also $\bits(R)\le2^{O(t)}L^{O(1)}$.

Define a deterministic answer rule $\mathcal O$ for integer queries $z$. If
$z\notin P_0$, return the first violated inequality of $P_0$ in a fixed order.
Otherwise compute $c_z$ by Lemma~\ref{lem:constraints-gradient-cut}. If
$c_z\ne0$, return $c_z^{\mathsf T}x\le c_z^{\mathsf T}z-\mu/4$. The answer
depends only on the instance and on $z$. The list algorithm starts with
$\mathcal Z=\{z_0\}$ and runs the algorithm of
Theorem~\ref{thm:constraints-integer-query} with radius $R$, answering
queries by $\mathcal O$ and adding every queried point of $P_0$ to
$\mathcal Z$. If some query $z\in P_0$ has $c_z=0$, it stops at once and
outputs $\mathcal Z$. Otherwise it outputs $\mathcal Z$ when the feasibility
algorithm halts.

\emph{Running time.} Let $\mathcal O_\varnothing$ agree with $\mathcal O$,
except that at a query $z\in P_0$ with $c_z=0$ it returns $x_1\le z_1-1$.
Every answer of $\mathcal O_\varnothing$ is an inequality violated by its
query, and every inequality is valid on the empty set. So
$\mathcal O_\varnothing$ is a valid oracle for $S=\varnothing$ that never
asserts membership, and its answers on queries of length $\ell$ cost at most
$\Phi(\ell)=(L+\bits(R)+\ell)^{C_1}$. Up to the possible early stop, the list
algorithm performs the run of Theorem~\ref{thm:constraints-integer-query} with
$\mathcal O_\varnothing$. That run halts, reporting emptiness, within
$2^{O(t\log(t+1))}\poly(t,1+\bits(R),\Phi(c_0t(1+\bits(R))))
=2^{O(t\log(t+1))}L^{C}$ steps, with $C$ absolute. The list and its total
length are bounded by the same quantity.

\emph{Every optimal block is listed.} If the algorithm stops early at $z$
with $c_z=0$, then $z$ is the unique optimal block by
Lemma~\ref{lem:constraints-integer-cut}, and it is listed. Otherwise suppose
an optimal block $w$ is not in $\mathcal Z$. Define the oracle $\mathcal O_w$
that asserts membership exactly at $w$ and answers every other query as
$\mathcal O$ does. It is a valid oracle for $S=\{w\}\subseteq[-R,R]^t$:
inequalities of $P_0$ are valid at $w\in P_0$, and at a query $z\in P_0$ with
$z\ne w$ we have $g(w)\le g(z)$, so $c_z\ne0$ and the cut is valid at $w$ by
Lemma~\ref{lem:constraints-integer-cut}. Its answer cost exceeds $\Phi$ only by
a comparison with the stored $w$. The actual run never queried $w$, and the
feasibility algorithm is deterministic, so the run with $\mathcal O_w$ asks
the same queries, receives the same answers, and also reports
$S\cap\Z^t=\varnothing$. This contradicts the correctness of the algorithm,
since $S\cap\Z^t=\{w\}$.

\emph{At most $2^t$ optimal blocks.} If $w_1\ne w_2$ are optimal blocks with
$w_1\equiv w_2\pmod 2$, then $(w_1+w_2)/2\in Q\cap\Z^t$ and strict convexity
of $g$ gives $g((w_1+w_2)/2)<\min g$. Hence distinct optimal blocks have
distinct parity vectors. For $f(z,y)=\sum_{i\le t}(z_i-\frac12)^2+\|y\|^2$ and
$Q=\R^t$, with $\mu=2$, the optimal blocks are exactly $\{0,1\}^t$.
\end{proof}

\begin{proof}[Proof of Corollary~\ref{cor:constraints-selection}]
If $Q\cap\Z^t=\varnothing$, the predicates have fixed truth values under the
convention $\min\varnothing=+\infty$. Otherwise the list $\mathcal Z$ contains
every optimal block, so $\min=\min_{z\in\mathcal Z}g(z)$. Thus
$\min\le\tau$ and $\min<\tau$ are disjunctions of the predicates
$g(z)\le\tau$ and $g(z)<\tau$, and the other relations are Boolean
combinations of these two. For $n\ge1$, the predicate $g(z)\bowtie'\tau$ is the
observable $f(z,y)-\tau$ at the minimizer of the strongly convex quartic
$y\mapsto f(z,y)$, so Theorem~\ref{thm:exact-upper} maps it to one $\PosSLP$
instance; for $n=0$ it is a rational comparison and maps to a fixed circuit.
All instances are built from $\mathcal Z$ without oracle calls, and the
Boolean form, Theorem~\ref{thm:posslp-closure}(c), combines them into one
instance. The total construction time is polynomial in the length of
$\mathcal Z$ and $L$, hence $2^{O(t\log(t+1))}L^{C}$.

For (b), compare every pair $z,z'\in\mathcal Z$ by
Corollary~\ref{cor:upper-two-minima}: the observable $f(z,y)-f(z',y')$ at the
minimizer of the strongly convex quartic $f(z,y)+f(z',y')$ on $\R^{2n}$, whose
Hessian is block diagonal and at least $\mu I$, equals $g(z)-g(z')$. A listed block is optimal exactly when its value
is at most every other listed value, because the list contains an optimal
block. All queries are fixed before any answer is received. The continuous
part of the optimum for an optimal block $z_*$ is $y(z_*)$, the unique real
solution of $\nabla_yf(z_*,y)=0$.
\end{proof}

\begin{lemma}[Padding that keeps a full Hessian Gram]
\label{lem:constraints-padding}
Let $f\in\Q[x_1,\ldots,x_n]$ be a quartic with a positive definite rational
full Hessian Gram $M$ on $(v,x\otimes v)$, of size $m_M=n+n^2$, and put
$\mu_0=\det M/(\tr M)^{m_M-1}$. For $t\ge1$ and
$\tau=1+8nt/\mu_0$, the quartic
\[
 P(x,z)=f(x)+\|z\|^2+\tau\|z\|^4+\|z\|^2\|x\|^2,\qquad z\in\R^t,
\]
has $\min\{P(x,z):x\in\R^n,\ z\in\Z^t\}=\min f$, attained exactly at
$z=0$, and a positive definite rational full Hessian Gram on the joint basis
$(v_x,v_z,x\otimes v_x,x\otimes v_z,z\otimes v_x,z\otimes v_z)$, computable in
polynomial time.
\end{lemma}

\begin{proof}
The added terms are nonnegative and vanish exactly when $z=0$, so the
mixed-integer minimum is $\min f$, attained only with $z=0$. In the direction
$(v_x,v_z)$, the second derivatives of the added terms are
$2\|v_z\|^2$, $\tau(4\|z\otimes v_z\|^2+8(z^{\mathsf T}v_z)^2)$, and
$2\|x\otimes v_z\|^2+2\|z\otimes v_x\|^2+8(x^{\mathsf T}v_x)(z^{\mathsf T}v_z)$.
Keep $M$ on $(v_x,x\otimes v_x)$; put $2I$ on $v_z$, $x\otimes v_z$ and
$z\otimes v_x$; put $4\tau I+8\tau\,\iota_t\iota_t^{\mathsf T}$ on $z\otimes v_z$,
where $\iota_t$ is the vectorized identity, so that
$\iota_t^{\mathsf T}(z\otimes v_z)=z^{\mathsf T}v_z$; and represent the
remaining term by the off-diagonal block $4\iota_n\iota_t^{\mathsf T}$ between
$x\otimes v_x$ and $z\otimes v_z$. For formal vectors $a$ in the
$x\otimes v_x$ block and $c$ in the $z\otimes v_z$ block, the cross term is at
most $8\sqrt{nt}\,\|a\|\|c\|\le\frac{\mu_0}2\|a\|^2+\frac{32nt}{\mu_0}\|c\|^2$
in absolute value. Since $M\succeq\mu_0I$
(Lemma~\ref{lem:models-det-trace}) and $4\tau-32nt/\mu_0=4$, the assembled
matrix is at least $\min\{\mu_0/2,2\}I$ on the whole basis, not only on
vectors of tensor form.
\end{proof}

\subsection{Mixed linear constraints}
\label{app:constraints-mixed}

Two further external results enter.

\begin{quote}
\emph{Contract (E3b)} \citep{KozlovTarasovKhachiyan1980}. A convex quadratic
function with rational coefficients that attains its minimum on a nonempty
explicit rational polyhedron has a rational minimizer that is computed exactly
in polynomial time.

\emph{Contract (E3e)} \citep[Theorem~3, applied with the zero
objective]{DelPia2023}. Given explicit rational $A$, $B$, $d$ with $t$ columns
of $A$, one decides in deterministic time $F_0(t)L^{C_0}$, with $F_0$
computable and $C_0$ absolute, whether
$\{(z,y):z\in\Z^t,\ Az+By\le d\}$ is empty, and otherwise returns a rational
point of it.
\end{quote}

\begin{lemma}[Primal--dual cuts for constrained fibers]
\label{lem:constraints-residual-cut}
Let $P=\{(z,y):Az+By\le d\}$, let $z\in\Z^t$ have a nonempty fiber, let
$\hat y$ be a rational point of the fiber and $\lambda\ge0$ a rational vector,
and put
\[
 \sigma=d-Az-B\hat y\ge0,\quad
 \rho=\nabla_yf(z,\hat y)+B^{\mathsf T}\lambda,\quad
 E=\lambda^{\mathsf T}\sigma+\frac{\|\rho\|^2}{2\mu},\quad
 c_z=\nabla_zf(z,\hat y)+A^{\mathsf T}\lambda .
\]
\begin{enumerate}[label=(\alph*)]
\item For every feasible $w$, $g(w)-g(z)\ge c_z^{\mathsf T}(w-z)+\frac\mu2\|w-z\|^2-E$.
In particular, if $E\le\mu/4$, then $c_z$ satisfies
\eqref{eq:constraints-integer-cut}.
\item If $n\ge1$ and $z\in[-R,R]^t$ for an integer $R\ge1$, then
$\hat y$ and $\lambda$ with $E\le\mu/4$ are computed in ordinary time
polynomial in $L+\bits(R)+\bits(z)$.
\end{enumerate}
\end{lemma}

\begin{proof}
(a) Let $(w,v)\in P$, and put $\Delta_z=w-z$ and $\Delta_y=v-\hat y$. Joint
strong convexity gives
\[
 f(w,v)\ge f(z,\hat y)+\nabla_zf^{\mathsf T}\Delta_z+\nabla_yf^{\mathsf T}\Delta_y
 +\tfrac\mu2\bigl(\|\Delta_z\|^2+\|\Delta_y\|^2\bigr),
\]
with derivatives at $(z,\hat y)$. Feasibility of $(w,v)$ gives
$A\Delta_z+B\Delta_y\le\sigma$, so
$\lambda^{\mathsf T}\sigma-\lambda^{\mathsf T}A\Delta_z-\lambda^{\mathsf T}B\Delta_y\ge0$.
Subtracting this quantity from the right side yields
\[
 f(w,v)\ge f(z,\hat y)+c_z^{\mathsf T}\Delta_z+\tfrac\mu2\|\Delta_z\|^2
 -\lambda^{\mathsf T}\sigma+\rho^{\mathsf T}\Delta_y+\tfrac\mu2\|\Delta_y\|^2,
\]
and $\rho^{\mathsf T}\Delta_y+\frac\mu2\|\Delta_y\|^2\ge-\|\rho\|^2/(2\mu)$.
Minimize over feasible $v$ and use $f(z,\hat y)\ge g(z)$. For integer $w\ne z$,
$\frac\mu2\|w-z\|^2-\frac\mu4\ge\frac\mu4\|w-z\|^2$.

(b) Write $h_z(y)=f(z,y)$ and let $y_z$ be its minimizer on the fiber
$F_z=\{y:By\le d-Az\}$. Linear programming gives a rational $\bar y\in F_z$,
and Lemma~\ref{lem:constraints-radius}(a), applied to $h_z$ and $\bar y$,
gives an integer $Y$ with $\|y_z\|<Y$. With $m_0=t+n$ and $C_f$ as before, put
\[
 S=2+R+Y,\quad \beta_1=1+4m_0C_fS^3,\quad \beta_2=1+12m_0C_fS^2,\quad
 \delta=\min\Bigl\{1,\frac{\mu}{16\beta_1},\frac{\mu}{4\beta_2}\Bigr\}.
\]
Then $\|\nabla h_z(y_z)\|\le\beta_1$, and $\|\nabla^2h_z(y)\|\le\beta_2$ when
$\|y-y_z\|\le1$, because all coordinates involved are at most $S$ in absolute
value. Contract (E1), applied to $h_z$ on $F_z\cap\{y:\|y\|_\infty\le Y\}$,
gives a rational $\hat y\in F_z$ with
$h_z(\hat y)\le h_z(y_z)+\min\{1,\mu\delta^2/2\}$, so
$\|\hat y-y_z\|\le\delta$ by Lemma~\ref{lem:constraints-radius}(b).

By Lemma~\ref{lem:constraints-polyhedral}(b), applied to $\nabla h_z$ on $F_z$,
there is a real $\lambda_*\ge0$, supported on the rows active at $y_z$, with
$\nabla h_z(y_z)+B^{\mathsf T}\lambda_*=0$. It need not be rational or
unique, and the proof uses no bound on its norm. Complementary slackness and
stationarity give the identity
\[
 \lambda_*^{\mathsf T}\sigma
 =\lambda_*^{\mathsf T}(d-Az-By_z)+(B^{\mathsf T}\lambda_*)^{\mathsf T}(y_z-\hat y)
 =\nabla h_z(y_z)^{\mathsf T}(\hat y-y_z),
\]
whose value lies in $[0,\beta_1\delta]$; it is nonnegative because $\hat y$ is
feasible and $y_z$ is optimal on $F_z$. Also
$\|\nabla h_z(\hat y)+B^{\mathsf T}\lambda_*\|=\|\nabla h_z(\hat y)-\nabla h_z(y_z)\|\le\beta_2\delta$.
Hence the convex quadratic function
\[
 \mathcal E(\lambda)=\sigma^{\mathsf T}\lambda
 +\frac{\|\nabla h_z(\hat y)+B^{\mathsf T}\lambda\|^2}{2\mu}
\]
satisfies
$\inf_{\lambda\ge0}\mathcal E\le\mathcal E(\lambda_*)\le\beta_1\delta+\beta_2^2\delta^2/(2\mu)\le\frac\mu{16}+\frac\mu{32}$.

The infimum is attained. The set
$\mathcal K=\{(B^{\mathsf T}\lambda,\sigma^{\mathsf T}\lambda):\lambda\ge0\}$ is a
finitely generated cone, hence closed, and its second coordinate is nonnegative.
In these coordinates the objective is $\tau+\|\nabla h_z(\hat y)+u\|^2/(2\mu)$,
whose sublevel set at the value of $\lambda=0$ is nonempty and compact within
$\mathcal K$. A minimizer $(u,\tau)\in\mathcal K$ lifts to some $\lambda\ge0$.
Contract (E3b) therefore returns an exact rational minimizer $\hat\lambda\ge0$,
and $E=\mathcal E(\hat\lambda)\le3\mu/32<\mu/4$. All data of the two convex
problems are rational of polynomial bit length, and the accuracy requested in
(E1) has polynomially many bits; no bound on $\|\lambda_*\|$ enters.
\end{proof}

\begin{proof}[Proof of Theorem~\ref{thm:constraints-mixed-candidates}]
Use contract (E3e) to certify infeasibility or to find a rational feasible
$(z_0,y_0)$. If $t=0$, output the empty block. Otherwise apply
Lemma~\ref{lem:constraints-radius}(a) to $f$ on $\R^{t+n}$ and $(z_0,y_0)$ to
get $R$ with every optimal block in $[-R,R]^t$; $\bits(R)\le F_1(t)L^{C_1}$.
Only the integer block is truncated; the fibers keep all their constraints.
Let $\mathcal D=\{z\in[-R,R]^t:\ \exists y,\ Az+By\le d\}$, a compact convex
set containing $z_0$ and every optimal block, on which $g$ is $\mu$-strongly
convex by the argument of Lemma~\ref{lem:constraints-gradient-cut}.

Answer an integer query $z$ as follows. If $z\notin[-R,R]^t$, return a
violated box inequality. If the fiber of $z$ is empty, linear programming
returns a rational $\alpha\ge0$ with $B^{\mathsf T}\alpha=0$ and
$\alpha^{\mathsf T}(d-Az)<0$; return $\alpha^{\mathsf T}Ax\le\alpha^{\mathsf T}d$,
which is violated by $z$ and valid on $\mathcal D$ because
$\alpha^{\mathsf T}(Az'+By')\le\alpha^{\mathsf T}d$ and
$\alpha^{\mathsf T}By'=0$. Otherwise record $z$, compute $c_z$ by
Lemma~\ref{lem:constraints-residual-cut} (or $c_z=\nabla f(z)$ if $n=0$), stop
if $c_z=0$, and return $c_z^{\mathsf T}x\le c_z^{\mathsf T}z-\mu/4$. These
answers depend only on the instance and the query, and cost polynomial time in
$L+\bits(R)+\bits(z)$.

The running-time argument (target $\varnothing$), the completeness argument
(target $\{w\}$ for an omitted optimal block $w$), and the early stop are
exactly as in the proof of Theorem~\ref{thm:constraints-candidates}: domain and
Farkas cuts are valid on $\mathcal D\ni w$, and residual cuts at $z\ne w$ are
valid at $w$ by Lemma~\ref{lem:constraints-integer-cut}, because
$g(w)\le g(z)$. The total time is $F_0(t)L^{C_0}$ plus
$2^{O(t\log(t+1))}\poly(\bits(R),L)$, which is $F(t)L^{C}$. For the parity
bound, if $(w_1,y_1)$ and $(w_2,y_2)$ are optimal with $w_1\ne w_2$ and
$w_1\equiv w_2\pmod 2$, their midpoint is feasible, has an integer block, and
has a strictly smaller objective value.
\end{proof}

\begin{proof}[Proof of Corollary~\ref{cor:constraints-nonlinear-mixed}]
Let $\mathcal Z$ be the list of Theorem~\ref{thm:constraints-mixed-candidates},
of total length at most $Q_L=F_1(t)L^{C_1}$. For $z\in\mathcal Z$, the fiber
instance $(f(z,\cdot),\{y:By\le d-Az\},\mu)$ has length at most
$(Q_L+L)^{C_2}$ with absolute $C_2$, because the degree is fixed. Its
curvature is at least $\mu$, since $\nabla^2_{yy}f$ is a principal block of
$\nabla^2f$. Write $f_3+f_4=\theta(U_zz+U_yy)$ with
$U=[U_z\ U_y]$ from Lemma~\ref{lem:constraints-essential}. For fixed $z$, the
homogeneous parts of degrees three and four in $y$ of $\theta(U_zz+U_yy)$ are
polynomials in $U_yy$, and $f_{\le2}(z,y)$ has degree at most two in $y$.
Hence the nonlinear kernel of $f(z,\cdot)$ contains $\ker U_y$, and its
nonlinear dimension is at most $k$. For $z,z'\in\mathcal Z$, the product
instance with objective $f(z,y)+f(z',y')$, feasible set $F_z\times F_{z'}$, and
observable $f(z,y)-f(z',y')$ has curvature at least $\mu$ and nonlinear
kernel containing $\ker U_y\times\ker U_y$, hence nonlinear dimension at most
$2k$. Theorem~\ref{thm:constraints-nonlinear} compares $g(z)$ and $g(z')$ in
time $F(2k)(Q_L+L)^{C_2C}$, and compares $g(z)$ with $\tau$ using
$h=f(z,y)-\tau$. At most $|\mathcal Z|^2+|\mathcal Z|$ comparisons select every
optimal block and decide the predicate, and
Theorem~\ref{thm:constraints-nonlinear}(iii) represents the continuous part.
The total time is $F'(k+t)L^{C'}$ with $F'(\kappa)$ the maximum of the
parameter factors over $k+t=\kappa$.
\end{proof}

\begin{proof}[Proof of Theorem~\ref{thm:constraints-mixed-rank}]
Compute the list $\mathcal Z$ of Theorem~\ref{thm:constraints-mixed-candidates},
remove duplicates, and sort it lexicographically. Each fiber has constraint
matrix $B$, of rank $r$, and curvature at least $\mu$. For $z,z'\in\mathcal Z$,
the product instance of the previous proof has constraint matrix
$\diag(B,B)$, of rank $2r$, so Theorem~\ref{thm:constraints-rank}, with
$T=\nabla(f(z,y)+f(z',y'))$, decides $g(z)<g(z')$ in expected time
$(2r+1)^{O(r)}(Q_L+L)^{C_3}$. Scan the sorted list, keeping the first block,
and replace the current block only by a block of strictly smaller value. The
result $z_*$ is the lexicographically smallest optimal block, because the list
contains every optimal block and ties keep the earlier one. One more call of
Theorem~\ref{thm:constraints-rank} on the fiber of $z_*$ returns a support $J$
with chart $(\bar y,Z)$ and the explicit strongly monotone cubic map $U_J$ whose
unique zero $u_*$ gives $y_*=\bar y+Zu_*$; the same call with
$h=f(z_*,y)-\tau$ decides $\min\bowtie\tau$. Comparing every listed block with
$z_*$ returns all optimal blocks.

The number of calls is at most $2|\mathcal Z|+2$, a deterministic bound.
Each call uses fresh random bits, and its expected running time, conditioned
on everything computed before it, is at most
$(2r+1)^{O(r)}(Q_L+L)^{C_3}$, because this bound holds for every input of
that length. Linearity of expectation bounds the total expected time by
$F(t,r)L^{C}$. Every call returns a correct answer, so every output is
correct. Different runs may return different charts, but all represent the
same point.
\end{proof}

\begin{proof}[Proof of Corollary~\ref{cor:constraints-mixed-unambiguous}]
Both machines first compute, deterministically, feasibility and the list
$\mathcal Z=\{z_1,\ldots,z_s\}$; infeasible instances are decided by the
convention. They then guess sets $S_1,\ldots,S_s$ of rows, one bit per row and
fiber, and verify each $S_i$ for the fiber of $z_i$ by steps (1)--(3) of the
proof of Theorem~\ref{thm:constraints-unambiguous}, with $T=\nabla_yf(z_i,\cdot)$.
By that proof, exactly one guess passes for every fiber, namely its active
set, and it yields a chart $(\bar y_i,Z_i)$ such that the fiber minimizer is
$\bar y_i+Z_iu_i$ with $u_i$ the unconstrained minimizer of
$q_i(u)=f(z_i,\bar y_i+Z_iu)$, whose Hessian is at least $\mu I$. The predicate
$\min\bowtie\tau$ is a Boolean combination of the predicates $q_i(u_i)\le\tau$
and $q_i(u_i)<\tau$, each decided by Theorem~\ref{thm:exact-upper} (or by
rational arithmetic for a zero-dimensional chart). One machine accepts when
the combination is true, the other when it is false. Each machine has at most
one accepting computation, and the guesses and verifications take time
polynomial in the length of the list.
\end{proof}

\subsection{A hard binary decision}
\label{app:constraints-binary}

\begin{proof}[Proof of Corollary~\ref{cor:constraints-binary}]
By Theorem~\ref{thm:rational-optimizer}, an integer straight-line program
with output $V$ is mapped in polynomial time to a rational quartic $F$ on
$\R^N$, an index $j$, $N+1$ rational quadratic square factors whose sum is
$F$, and a positive definite rational full Hessian Gram $Q$, such that $F\ge0$
has a unique zero $p\in\Q^N\cap[-1,1]^N$, $\nabla^2F\succeq\frac32I$,
$p_j\ne0$, and $p_j>0$ exactly when $V>0$.

In \eqref{eq:constraints-binary}, the integer $z$ is $0$ or $1$, and then
$H=F(x)+\frac14\ge\frac14$, with equality only at $x=p$. The constraint
$z-1\le x_j\le z$ allows $x_j\in[-1,0]$ for $z=0$ and $x_j\in[0,1]$ for
$z=1$. Both fibers are nonempty ($x=0$), and since $0<|p_j|\le1$, exactly the
fiber $z=\mathbf 1[p_j>0]$ contains $p$. Hence the optimal value is $1/4$,
the unique optimal solution is $(p,\mathbf 1[p_j>0])$, and $z_*=1$ exactly
when $V>0$. The Hessian of $H$ is $\diag(\nabla^2F,2)\succeq\frac32I$. The
constraints $0\le z\le1$ have zero continuous part, and $z-1\le x_j\le z$
have continuous normals $\pm e_j$, so the continuous constraint rank is one.
The other coordinates of $x$ are unconstrained; the optimizer is bounded,
although the feasible set is not.

For the full Gram, let $\varrho=\det(Q)/\tr(Q)^{d_Q-1}$, where $d_Q=N+N^2$ is
the size of $Q$, so $Q\succeq\varrho I$ (Lemma~\ref{lem:models-det-trace}),
and let $\epsilon=\min\{\varrho/4,1/(100N)\}$. Put
\[
 \widehat H=H+\epsilon z(z-1)\|x\|^2+\tfrac14z^2(z-1)^2 .
\]
Both added terms vanish at $z\in\{0,1\}$, so every feasible objective value,
the optimum, the optimizer, and the constraints are unchanged. Write
$\vartheta=z-\frac12$, let $(v,s)$ be a direction in $(x,z)$, and let
$w=(v,x\otimes v)$. Direct differentiation of
$\widehat H=F+\vartheta^2+\epsilon(\vartheta^2-\frac14)\|x\|^2+\frac14(\vartheta^2-\frac14)^2$
gives
\[
 D^2\widehat H[(v,s)]^2=w^{\mathsf T}Qw-\tfrac\epsilon2\|v\|^2+\tfrac74s^2
 +2\epsilon\|\vartheta v\|^2+2\epsilon\|xs\|^2+8\epsilon(x^{\mathsf T}v)(\vartheta s)
 +3(\vartheta s)^2 .
\]
In the basis $(w,s,\vartheta v,xs,\vartheta s)$, which contains every
direction coordinate and every product of a coordinate of $(x,z)$ with a
direction coordinate exactly once after the invertible rational change
$\vartheta v=zv-\frac12v$, $\vartheta s=zs-\frac12s$, this is the quadratic form
of the block matrix with diagonal blocks $Q-\frac\epsilon2E$, $\frac74$,
$2\epsilon I$, $2\epsilon I$, $3$, and one off-diagonal block $4\epsilon d_0$
between $w$ and $\vartheta s$. Here $E$ is the identity on the $v$-part of $w$
and zero elsewhere, and $d_0$ is the vector with $d_0^{\mathsf T}w=x^{\mathsf T}v$,
so $\|d_0\|^2=N$. Since $Q-\frac\epsilon2E\succeq\frac\varrho2I$, the Schur
complement of the last entry is at least
$3-32\epsilon^2N/\varrho\ge3-8\epsilon N\ge3-\frac2{25}>0$. Hence the matrix
is positive definite, and so is its rational congruence in the original basis.
All entries have polynomial bit length.
\end{proof}

The base objective $H$ has no positive definite full joint Hessian Gram: at
$x=0$ in the direction $s$, its Hessian form is the constant $2s^2$, while a
positive definite Gram on a basis containing $zs$ would force growth like
$z^2s^2$. The correction restores the full Gram at the cost of the specific
modulus $3/2$ and of the supplied objective square factors. The value gap
between the two fibers is at least $\frac34p_j^2$, by strong convexity and the
distance from $p$ to the wrong half-space, and at most $\frac{\Lambda_F}2p_j^2$
with $\Lambda_F=\max_{[-1,1]^N}\|\nabla^2F\|$, by evaluating at $p-p_je_j$.
